\documentclass[11pt]{article}

\usepackage[margin=1in]{geometry}
\usepackage{amsmath,amssymb,amsthm,mathtools}
\usepackage{booktabs,array,longtable,microtype}
\usepackage[hidelinks]{hyperref}
\IfFileExists{cleveref.sty}{%
  \usepackage[nameinlink,noabbrev]{cleveref}%
}{%
  \newcommand{\cref}[1]{\autoref{##1}}%
}

\hypersetup{
  pdftitle={Separator compression of the complete spin-structure family for cubic-lattice Ising strips},
  pdfauthor={Hainan Zhao}
}

\newtheorem{theorem}{Theorem}[section]
\newtheorem{proposition}[theorem]{Proposition}
\newtheorem{lemma}[theorem]{Lemma}
\newtheorem{corollary}[theorem]{Corollary}
\newtheorem{definition}[theorem]{Definition}
\newtheorem{remark}[theorem]{Remark}

\newcommand{\Ftwo}{\mathbb F_2}
\newcommand{\cE}{\mathcal E}
\newcommand{\cF}{\mathcal F}
\newcommand{\cQ}{\mathcal Q}
\newcommand{\rank}{\operatorname{rank}}
\newcommand{\Flat}{\operatorname{Flat}}
\newcommand{\Arf}{\operatorname{Arf}}
\newcommand{\im}{\operatorname{im}}
\newcommand{\supp}{\operatorname{supp}}
\newcommand{\Zar}{\operatorname{Zar}}
\newcommand{\claimtag}[1]{\textsc{#1}}
\newcommand{\proved}{\claimtag{Proved}}
\newcommand{\certified}{\claimtag{Certified\_Numerical}}
\newcommand{\observed}{\claimtag{Observed}}
\newcommand{\conjectured}{\claimtag{Conjectured}}
\newcommand{\proofpending}{\claimtag{Proof incomplete}}

\title{Separator Compression of the Complete Spin-Structure Family\\
for Cubic-Lattice Ising Strips}
\author{Hainan Zhao}
\date{August 2026\\\small Reserved archive DOI: \url{https://doi.org/10.5281/zenodo.21845273}}

\begin{document}
\maketitle

\begin{abstract}
For a graph cellularly embedded in an orientable surface of genus \(g\), the
generalized Kac--Ward formula represents the zero-field Ising even-subgraph
polynomial as an Arf-weighted sum of \(4^g\) signed square roots.  We study the
complete tensor of those pre-Arf values, before the final sum.  We prove an
abstract separator theorem: if a filtration-adapted spin-structure basis makes every
cross-cut linear and quadratic phase factor through the even boundary mask on
a separator \(S\), then every handle-pair flattening has rank at most
\(2^{|S|-1}\).  A relative exactness condition gives the same bound at cuts
inside a handle.

For the cubic grid strips

\[
  G_{n,w}=P_n\mathbin\square P_w\mathbin\square P_w,
\]

we give an explicit nested checkerboard ribbon embedding and a candidate
canonical handle ordering with \(|S|=w^2\).  The target conclusion is that
the entire \(4^g\)-entry spin-structure family has fraction-field bond at
most \(d_w=2^{w^2-1}\), and that for independent nonuniform weights and every
\(w\geq3\), some canonical cut generically attains this bound with \(n=11\)
uniformly sufficient.  The remaining self-contained proof obligations are
marked in the text.  If discharged, these statements show that surface
genus creates no additional exact bond dimension beyond the ordinary
zero-field transverse carrier, but that carrier is itself generically
minimal.  We also give a non-grid toroidal \(K_{3,3}\) chain showing that the
internal-handle hypothesis is sharp, and an \(O(gd_w^2)\) contraction for all
single-handle Walsh marginals of the complete family.

This is a structural and exact-complexity result for the spin-structure
representation.  For cubic boxes \(w=n=L\), the remaining bond is
\(2^{L^2-1}\).  No thermodynamic limit, critical temperature, or exact solution
of the three-dimensional Ising model is obtained.
\end{abstract}

\begin{center}
\fbox{\begin{minipage}{.92\textwidth}\small
\textbf{Revision status.}
The headline grid upper bound and generic-saturation statements are retained
as targets but are marked \proofpending.  Exact finite-field computations in
this draft are audits, not substitutes for the arbitrary-width geometric and
combinatorial proofs.  In particular, no submission-ready claim is made
until the explicitly identified proof obligations are discharged.
\end{minipage}}
\end{center}

\section{Introduction}

Let \(G=(V,E)\) be a finite graph with zero magnetic field and couplings
\(K_e\).  With \(t_e=\tanh K_e\), the high-temperature expansion is
\begin{equation}\label{eq:ht}
 Z_G=2^{|V|}\prod_{e\in E}\cosh K_e\;\cE_G(t),
 \qquad
 \cE_G(t)=\sum_{A\in Z_1(G;\Ftwo)}\prod_{e\in A}t_e .
\end{equation}
Here \(Z_1(G;\Ftwo)\) is the set of even edge subsets.  If \(G\) is embedded
in an orientable surface of genus \(g\), the generalized Kac--Ward formula
expresses \(\cE_G\) as an Arf-weighted sum over all \(4^g\) quadratic
refinements of the surface intersection form
\cite[Theorem~2.1]{Cimasoni2010}.  Equivalently, these are the signed
Kac--Ward square roots or Fisher--Kasteleyn Pfaffians associated with the
surface spin structures.

The number of terms is topologically natural and, within the corresponding
linear-combination model, optimal: Loebl and Masbaum prove that \(4^g\)
projected Pfaffians are necessary at embedding genus \(g\)
\cite[Theorems~4 and 4.1]{LoeblMasbaum2011}.  Our question is different.
Rather than replace the \(4^g\) functions by fewer Pfaffians, we preserve all
of them and ask for the exact tensor rank of their joint dependence on the
spin-structure coordinates.  This distinction is essential: none of the
results below contradicts Pfaffian-count optimality.

For fixed \(w\), the strip \(G_{n,w}\) has bounded pathwidth and is already
tractable by ordinary transfer or tree-decomposition dynamic programming
\cite{ScottSorkin2009}.  The issue is therefore not whether a fixed-width
strip can be evaluated.  It is whether the exponentially many surface
sectors carry an additional topological memory beyond the ordinary
transverse state.  The target theorems below would answer that question
exactly once their marked proof obligations are discharged.

\subsection{Main results and epistemic boundary}

Completed mathematical arguments are marked \proved, exact finite-field
cross-checks are marked \certified, and claims whose self-contained proof is
still being assembled are marked \proofpending.

For fixed \(w\), write
\begin{equation}\label{eq:Rinf}
 R_\infty(w)=\sup_n\max_{\text{canonical pair and internal cuts}}
       \rank_{\operatorname{Frac}\mathbb Z[t_e]}\Flat(\cF_{n,w}).
\end{equation}
This is the generic, or fraction-field, rank invariant.  Equality means that
at least one canonical cut attains the carrier bound; it does not assert that
every cut or every specialization has that rank.

\begin{enumerate}
\item \proved: The abstract separator theorem, \cref{thm:separator}, bounds
  the complete pre-Arf tensor by the attainable separator-mask space.  Its
  hypotheses are intrinsic relative-chain exactness and quadratic gluing,
  not planarity, minimum genus, or translation invariance.

\item \proofpending: For an explicit checkerboard embedding of \(G_{n,w}\), every
  pair cut and every binary cut inside a canonical handle has rank at most
  \[
     d_w=2^{w^2-1}.
  \]
  Consequently the complete function of all \(2g\) spin-structure bits has
  a fraction-field tensor train of that bond dimension
  (\cref{thm:grid-upper}).  The denominator-free polynomial-ring refinement
  remains a separate proof obligation.

\item \proofpending: For generic independent nonuniform edge weights and every
  \(w\geq3\), the preceding bound is minimal:
  \begin{equation}\label{eq:g1-intro}
     R_\infty(w)=2^{w^2-1}.
  \end{equation}
  A length \(n=11\) suffices uniformly in \(w\)
  (\cref{thm:grid-tight}).  The lower bound holds on a nonempty open set of
  strictly positive ferromagnetic weights \(0<t_e<1\), but it is not proved
  after restricting to homogeneous
  anisotropic or isotropic weights.

\item \proofpending: An iterated toroidal \(K_{3,3}\) two-sum chain has generic
  pair rank two but generic internal-handle rank four.  This is an infinite
  sharpness family for the internal relative-exactness condition
  (\cref{thm:k33}).

\item \proofpending: Given a denominator-free all-sector tensor train, all four
  single-handle Walsh marginals at every handle, under arbitrary product-form
  sector weights, are obtained in \(O(gd_w^2)\) ring operations
  (\cref{thm:marginals}).
\end{enumerate}

The conclusion is deliberately limited.  The surface-spin-structure factor
has no rank overhead beyond the physical separator carrier.  For cubic boxes
\(w=L\), however, that carrier has dimension \(2^{L^2-1}\), and
\cref{eq:g1-intro} shows it is generically unavoidable within this complete
pre-Arf tensor.  We do not derive a controlled thermodynamic limit, an exact
free energy, a critical temperature, or critical exponents.

\subsection{Relation to existing mechanisms}

Cimasoni's theorem extends the planar Kac--Ward mechanism
\cite{KacWard1952} to the surface Kac--Ward/Arf representation, but does
not organize the complete family as a separator-adapted tensor train
\cite{Cimasoni2010}.  Bravyi contracts nonplanar matchgate networks by
Grassmann Gaussian integration with an explicit genus-dependent sector sum
\cite{Bravyi2009}; our result instead isolates when all sector phases factor
through the same physical boundary mask.  Gobron packages the nonplanar
calculation into one Pfaffian over a multicomplex coefficient algebra whose
size still grows with the embedding parameter \cite[Theorem~3.2 and
Corollary~3.4]{Gobron2013}; our rank statement concerns the collective sector
tensor over the ordinary edge-weight ring.  Standard bounded-width dynamic
programming supplies a \(2^{\Theta(w^2)}\) physical frontier but does not by
itself state the all-spin-structure factorization.  The theorem below is the
bridge between these two descriptions.

The minimum genus of many grid products has been studied by Millichap and
Salinas \cite{MillichapSalinas2022}.  We use their Theorem~4 only for the
separate statement that \(P_n\square P_3\square P_3\) has minimum orientable
genus \(n-1\).  Minimum genus is not a hypothesis of our compression theorem.
No broad priority claim is made beyond the theorem-by-theorem comparisons
just stated.

\section{The complete pre-Arf tensor}

Let \(G\hookrightarrow\Sigma_g\) be a cellular embedding in a closed
orientable surface and let
\[
 \pi:Z_1(G;\Ftwo)\longrightarrow H_1(\Sigma_g;\Ftwo)=H
\]
be induced by inclusion.  Fix a symplectic basis
\((a_1,b_1,\ldots,a_g,b_g)\), so the intersection form is
\(J=\bigoplus_{i=1}^g\left(\begin{smallmatrix}0&1\\1&0\end{smallmatrix}\right)\).
For
\(\lambda=(\lambda_{a_1},\lambda_{b_1},\ldots,
\lambda_{a_g},\lambda_{b_g})\in\Ftwo^{2g}\), define
\begin{equation}\label{eq:q}
 q_\lambda(x,y)=x\cdot y+\lambda_a\cdot x+\lambda_b\cdot y,
 \qquad (x,y)\in\Ftwo^g\oplus\Ftwo^g.
\end{equation}
Then
\[
 q_\lambda(u+v)=q_\lambda(u)+q_\lambda(v)+u^TJv,
 \qquad
 \Arf(q_\lambda)=\lambda_a\cdot\lambda_b.
\]

\begin{definition}[complete pre-Arf tensor]\label{def:F}
Over the polynomial ring \(R=\mathbb Z[t_e:e\in E]\), put
\begin{equation}\label{eq:F}
 \cF_G(\lambda)
  =\sum_{A\in Z_1(G;\Ftwo)}
       (-1)^{q_\lambda(\pi A)}\prod_{e\in A}t_e.
\end{equation}
The order-\(g\), four-state tensor obtained by grouping
\((\lambda_{a_i},\lambda_{b_i})\) is the complete pre-Arf tensor.
\end{definition}

By character orthogonality,
\begin{equation}\label{eq:arf}
 \cE_G(t)=2^{-g}\sum_{\lambda\in\Ftwo^{2g}}
     (-1)^{\lambda_a\cdot\lambda_b}\cF_G(\lambda).
\end{equation}
This is also the signed square-root form of the generalized Kac--Ward
identity, after fixing its branches and global signs consistently
\cite[Theorem~2.1]{Cimasoni2010}.  We use \cref{eq:F}, which is polynomial
and branch-free, as the definition; the determinant representation supplies
the interpretation, not an unverified identity.

For an ordered tensor \(T(s_1,\ldots,s_g)\), the rank at the cut
\(\{1,\ldots,j\}\mid\{j+1,\ldots,g\}\) is the matrix rank of its flattening over
the fraction field of \(R\).  A tensor train of bond \(D\) is an identity
\begin{equation}\label{eq:tt}
 T(s_1,\ldots,s_g)=\ell^TA_1(s_1)\cdots A_g(s_g)r
\end{equation}
with all intermediate dimensions at most \(D\).  We also use the binary
ordering
\((\lambda_{a_1},\lambda_{b_1},\ldots,\lambda_{a_g},\lambda_{b_g})\), which
adds cuts inside a handle.

\section{An abstract separator theorem}

We formulate the mechanism independently of grids.  Partition
\(E=E_L\sqcup E_R\) across a labelled interface \(S\).  Partial chains have
boundary maps
\[
 \partial_L:\Ftwo^{E_L}\to\Ftwo^S,
 \qquad
 \partial_R:\Ftwo^{E_R}\to\Ftwo^S,
\]
and glue to an even subgraph exactly when their boundary masks agree.  For a
connected interface, every attainable mask lies in
\begin{equation}\label{eq:VS}
 V(S)=\left\{m\in\Ftwo^S:\sum_{s\in S}m_s=0\right\},
 \qquad |V(S)|=2^{|S|-1}.
\end{equation}
Choose linear completion chains \(c_L(m),c_R(m)\), and write
\(h_L(A_L,m),h_R(A_R,m)\) for the homology classes obtained by closing a
partial chain.

\begin{definition}[filtration-adapted cut]\label{def:adapted}
A cut of the ordered spin coordinates into \(\lambda_L|\lambda_R\) is
\emph{filtration adapted} if the following conditions hold.

\begin{description}
\item[H1: relative character exactness.]
Every character cochain belonging to \(\lambda_L\), restricted to right
partial chains, is a coboundary plus a fixed linear function of the trace
\(m\); the left-right reversed assertion holds for \(\lambda_R\).
Concretely, put
\[
 P_R=\{c\in C_1(G_R;\Ftwo):\supp(\partial c)\subseteq S\},
 \qquad \tau_R(c)=\partial c|_S,
\]
and define \(P_L\) analogously.  There are
\(s\in C^0(G_R;\Ftwo)\) and a linear
functional \(\phi\) on \(\im\tau_R\) such that
\begin{equation}\label{eq:H1}
 \langle\alpha,c\rangle
   =\langle\delta s,c\rangle+\phi(\tau_Rc)
   =\langle s,\tau_Rc\rangle+\phi(\tau_Rc)
\end{equation}
for every \(c\in P_R\).

\item[H2: quadratic gluing.]
There are binary functions \(Q_L,Q_R,\kappa\) such that for all compatible
partial chains,
\begin{equation}\label{eq:H2}
 q_0(\pi(A_L+A_R))
   =Q_L(A_L,m)+Q_R(A_R,m)+\kappa(m).
\end{equation}

\item[H3: internal-handle trace.]
At a cut \(\lambda_{a_i}|\lambda_{b_i}\), the evaluation of the un-emitted
\(b_i\)-character, represented by the cochain \(\beta_i\), on the exposed
half depends only on its trace: if \(c,c'\)
are exposed partial chains with \(\tau(c)=\tau(c')=m\), then
\(\langle\beta_i,c\rangle=\langle\beta_i,c'\rangle\).  Consequently
\begin{equation}\label{eq:rho}
 \rho_i(m):=\langle\beta_i,c\rangle
 \quad\text{for any exposed }c\text{ with }\tau(c)=m
\end{equation}
is well defined.
\end{description}
\end{definition}

Affine changes of quadratic origin are linear characters and are included in
H1.  Literal cochain exactness is sufficient but not necessary; what is
intrinsic is equality of the induced phase on partial configurations with the
same separator trace.

\begin{theorem}[separator compression; \proved]\label{thm:separator}
Suppose every pair cut in a canonical handle ordering has a cut datum with
\(|S|\leq k\) satisfying H1--H2.  Then every pair flattening of the complete
pre-Arf tensor has rank at most \(2^{k-1}\).  If H3 holds at every internal
handle cut, the same bound holds in the binary-coordinate ordering.  Without
H3, the general binary-coordinate bound is \(2^k\).
\end{theorem}

\begin{proof}
Fix a cut.  By H1, every left spin coordinate evaluated on a right partial
chain is a sign depending only on \(m\), and conversely for right
coordinates.  Absorb these signs into the corresponding factors.  H2 then
gives, coefficientwise in the independent edge variables,
\begin{equation}\label{eq:XY}
 \cF_G(\lambda_L,\lambda_R)
 =\sum_{m\in V(S)}X(\lambda_L,m)(-1)^{\kappa(m)}
                       Y(m,\lambda_R),
\end{equation}
where, for example,
\[
 X(\lambda_L,m)=
 \sum_{\substack{A_L\subseteq E_L\\\partial_LA_L=m}}
 (-1)^{Q_L(A_L,m)+\text{left character phases}}
 \prod_{e\in A_L}t_e .
\]
Thus the flattening factors through the attainable mask space, of cardinality
at most \(2^{k-1}\).

At an internal handle cut, the pre-Arf sign is still the sign in
\cref{eq:q}: its dependence on \(\lambda_{b_i}\) is the character factor
\((-1)^{\lambda_{b_i}\langle\beta_i,c\rangle}\).  By H3 and
\cref{eq:rho}, this is the mask diagonal
\((-1)^{\lambda_{b_i}\rho_i(m)}\), so it can be absorbed into the factor
containing \(\lambda_{b_i}\) without enlarging the shared index.  The
quadratic term \(x_i y_i\) belongs to \(q_0\) and has already been split by
H2.  (The product \(\lambda_{a_i}\lambda_{b_i}\) occurs in the final Arf
weight, not in the pre-Arf tensor.)  Without H3,
splitting a four-state core into its two binary coordinates gives the general
factor-two bound.  The bounds at all successive cuts imply a tensor train
over the fraction field of \(R\) by the successive-rank characterization of
tensor-train rank.  This implication alone does not produce
denominator-free cores over \(R\).
\end{proof}

For a disconnected interface, the correct bound is the cardinality of the
actually attainable trace module; it need not equal \(2^{|S|-1}\).
Free boundaries select the zero mask.  An antiperiodic seam is a fixed edge
cocycle and changes only a local sign.  A periodic closure is a trace over the
same mask carrier after cutting a seam.  Fixed-spin boundaries must be
transformed together with a gauge change.

\section{The checkerboard grid embedding}

Put
\[
 V_{n,w}=\{(i,y,z):0\leq i<n,\ 0\leq y,z<w\},
\]
with nearest-neighbour grid edges, and decompose along the \(i\)-coordinate.
The transverse slice is
\(\Lambda_w=P_w\square P_w\).

\subsection{Rotation system and genus}

Let \(N,W\) be the least even integers with \(N\geq n\), \(W\geq w\).
For \(c=(c_0,c_1,c_2)\) with
\(0\le c_0<N-1\) and \(0\le c_1,c_2<W-1\), put
\begin{equation}\label{eq:cube-set}
 c\in\mathcal C_{N,W}\quad\Longleftrightarrow\quad
 \#\{r:c_r\equiv0\pmod2\}\ge2.
\end{equation}
This parity rule determines the rotation system as follows.  If the face of
the cube \(c\) perpendicular to axis \(a\), with outward sign
\(\eta\in\{-1,+1\}\), is not shared with another cube in
\(\mathcal C_{N,W}\), orient its four vertices
\(v_0,v_1,v_2,v_3\) by the outward-normal rule and set
\begin{equation}\label{eq:successor}
 \operatorname{succ}_{v_r}(v_{r-1})=v_{r+1}\qquad(r\bmod4).
\end{equation}
For reference, if \(b<c\) are the two axes different from \(a\), the positive
ordered tangent pair is
\begin{center}
\begin{tabular}{c|rrrrrr}
\((a,\eta)\)&\((0,+)\)&\((0,-)\)&\((1,+)\)&\((1,-)\)&\((2,+)\)&\((2,-)\)\\
\midrule
\((b,c)\)&\((1,2)\)&\((2,1)\)&\((2,0)\)&\((0,2)\)&\((0,1)\)&\((1,0)\)
\end{tabular}
\end{center}
and \(v_0,v_1,v_2,v_3\) traverses the corresponding unit square in the
positive \((b,c)\) order.

Here is the parity check that makes the definition global.  Around an
interior edge, the four adjacent cube indicators, read cyclically, are a
cyclic shift of \(1110\) when the edge coordinate is even and of \(1000\)
when it is odd.  In either case there are exactly two indicator changes.
Thus exactly two oriented boundary faces contain every grid edge, and they
give one predecessor and one successor at each endpoint.  The vertex link is
therefore a single oriented cycle; its neighbour order is the cycle of
\(u\mapsto\operatorname{succ}_v(u)\).  At a boundary vertex, out-of-box cube
indicators are zero.  Equivalently, its cyclic order is obtained by deleting
absent neighbours from the next-even parent box.  This covers all eight
vertex parities and every boundary type without an implicit geometric
choice.

Let \(D\) be the dart set of the retained grid, let \(\alpha\) reverse a
dart, and let \(\sigma\) advance it in the preceding vertex order.  The
orbits of \(\varphi=\sigma\alpha\) are the face walks.  Attach one oriented
disk along each orbit.  Every edge band then has two oppositely oriented
sides and every vertex link is one circle, so the resulting CW complex is a
closed orientable surface.  Its complement is exactly the union of the
attached open disks, proving cellularity directly.  We denote this surface
by \(\Sigma_{n,w}\).  Deleting a terminal longitudinal slice deletes
precisely its darts from \(\sigma\), so the restricted rotations are nested
in \(n\).

\begin{figure}[t]
\centering
\setlength{\unitlength}{.7mm}
\begin{picture}(115,45)
 \thicklines
 \put(5,5){\line(1,0){35}}\put(40,5){\line(0,1){28}}
 \put(40,33){\line(-1,0){35}}\put(5,33){\line(0,-1){28}}
 \put(5,33){\line(2,1){18}}\put(40,33){\line(2,1){18}}
 \put(40,5){\line(2,1){18}}\put(58,23){\line(0,1){19}}
 \put(23,42){\line(1,0){35}}
 \put(15,18){selected cube}
 \put(8,8){\vector(1,0){23}}\put(37,9){\vector(0,1){18}}
 \thinlines
 \put(62,23){\vector(1,0){13}}
 \put(78,19){\shortstack{oriented boundary faces\\define the dart successor}}
\end{picture}
\caption{Local ribbon rule.  A boundary plaquette is oriented by its outward
normal; \eqref{eq:successor} converts its two incident sides at a vertex into
the cyclic-order successor.}
\label{fig:rotation}
\end{figure}

\begin{proposition}[constructed genus; \proofpending]\label{prop:genus}
Let \(p\geq1\).  The genus of the preceding embedding is
\begin{align}
 g(n,2)&=0,\nonumber\\
 g(n,2p+1)&=p^2(n-1),\label{eq:genus}\\
 g(n,2p)&=\sum_{i=0}^{n-2}r_i,
 &r_i&=\begin{cases}(p-1)^2,&i\text{ even},\\p^2-1,&i\text{ odd}.
 \end{cases}\nonumber
\end{align}
When \(n,w\) are even the embedding is minimum genus.  For \(w=3\), its
genus \(n-1\) is minimum by Millichap--Salinas, Theorem~4
\cite{MillichapSalinas2022}.  No other minimum-genus assertion is needed.
\end{proposition}

\begin{proof}
We trace the face permutation \(\varphi=\sigma\alpha\).  For the first slice,
the ribbon graph is the planar \(w\times w\) grid, so
\[
 F(1,w)=2w(w-1)-w^2+2=w^2-2w+2.
\]
When slab \(i\) is attached, the new vertices and edges number
\[
 \Delta V=w^2,\qquad
 \Delta E=w^2+2w(w-1)=3w^2-2w.
\]
The \(\varphi\)-orbits meeting the slab can be followed from their first
transverse dart.  The parity rule \eqref{eq:cube-set} gives the following
exhaustive trace table:
\[
\begin{array}{c|c|c}
 w& i\pmod2&r_i\\ \hline
 2p+1&0,1&p^2\\
 2p&0&(p-1)^2\\
 2p&1&p^2-1.
\end{array}
\qquad
\Delta F=2w(w-1)-2r_i.                              \tag{*}
\]
To verify the table, scan the transverse squares by their lower corner
\((y,z)\).  A face orbit entering at that square exits at the translated
square in the old boundary unless
\[
 y\equiv z\equiv i+1\pmod2.
\]
At such a square its four local darts instead split the old orbit into the
two boundary components of a new band.  There are \(p^2\) such squares at
odd width.  At even width there are \((p-1)^2\) for \(i\) even; for \(i\)
odd there are \(p^2\), but their outer arcs concatenate into one orbit, so
the sum of all \(p^2\) local splittings has one relation.  Every transverse
square belongs either to this congruence class or to a translated old orbit,
and the longitudinal boundary rows supply the remaining
\(2w(w-1)-2\#\{\text{independent splittings}\}\) face orbits.  This proves
both the exhaustion and \((*)\) directly from \(\varphi\).

Euler's formula now gives
\[
 \Delta(2-2g)=\Delta V-\Delta E+\Delta F=-2r_i.
\]
The first slice has genus zero, so summing the displayed \(r_i\) proves
\cref{eq:genus}.

For a numerical check on the face count, at \(w=3\) the rotation system has
\[
 V=9n,\qquad E=21n-9,\qquad F=10n-5,
\]
and hence \(V-E+F=4-2n=2-2(n-1)\).  If \(n,w\) are even, every face is a
quadrilateral.  Since
\[
 V=nw^2,\qquad E=(n-1)w^2+2nw(w-1),\qquad F=E/2,
\]
the constructed genus is
\begin{equation}\label{eq:even-genus-euler}
 g=1+\frac{E-2V}{4}.
\end{equation}
In any cellular embedding of this bipartite graph, girth four gives
\(4F\leq2E\); Euler's formula therefore gives
\(g\geq1+(E-2V)/4\).  Equality in \cref{eq:even-genus-euler} proves the
minimum-genus assertion.
\end{proof}

\begin{figure}[t]
\centering
\setlength{\unitlength}{.65mm}
\begin{picture}(125,48)
 \multiput(5,8)(10,0){5}{\line(0,1){40}}
 \multiput(5,8)(0,10){5}{\line(1,0){40}}
 \multiput(78,8)(10,0){5}{\line(0,1){40}}
 \multiput(78,8)(0,10){5}{\line(1,0){40}}
 \thicklines
 \put(5,8){\framebox(10,10){}}\put(25,8){\framebox(10,10){}}
 \put(5,28){\framebox(10,10){}}\put(25,28){\framebox(10,10){}}
 \put(78,8){\framebox(10,10){}}\put(98,8){\framebox(10,10){}}
 \put(78,28){\framebox(10,10){}}\put(98,28){\framebox(10,10){}}
 \put(52,27){\vector(1,0){18}}\put(51,31){attach slab}
 \put(82,13){\oval(12,7)}\put(88,13){\vector(0,1){22}}
 \put(12,2){co-core meridians}\put(88,2){dual longitude}
\end{picture}
\caption{One slab attachment.  The congruence-class transverse squares are
co-core meridians; cutting them leaves the translated old face orbits and a
planar collar.  A longitude crosses its meridian once and the others zero
times.}
\label{fig:slab-cocores}
\end{figure}

\subsection{A nested canonical Lagrangian}

Let \(P_{i,y,z}\) be the four-edge boundary of the transverse unit square in
slice \(i\) with lower corner \((i,y,z)\).  In the slab from \(i\) to \(i+1\),
take the classes
\begin{equation}\label{eq:Ais}
 [P_{i,y,z}],\qquad y\equiv z\equiv i+1\pmod2.
\end{equation}
If \(w=2p\) and \(i\) is odd, omit the lexicographically last of these \(p^2\)
classes.  Order first by slab and then lexicographically.

\begin{lemma}[canonical handle basis; \proofpending]\label{lem:canonical}
The classes in \cref{eq:Ais} form a Lagrangian basis
\(A=(a_1,\ldots,a_g)\) of \(H_1(\Sigma_{n,w};\Ftwo)\).  They admit a nested
symplectic extension \(B=(b_1,\ldots,b_g)\), and the resulting order
\((a_1,b_1,\ldots,a_g,b_g)\) is filtration adapted.
\end{lemma}

\begin{proof}
Within a slab the squares are disjoint; different slabs use different
transverse slices.  Hence their mutual intersections vanish.  They are the
boundaries of the co-core disks shown in \cref{fig:slab-cocores}.  Completeness
is a consequence of the face trace, not merely of the count.  In one slab,
cut along all congruence-class disks and contract each translated old face
orbit and each longitudinal boundary strip.  The remaining incidence graph
has a vertex for each boundary row and an edge for each non-co-core square.
Scanning first in \(y\) and then in \(z\) gives the path
\[
 (0,0),(1,0),\ldots,(w-1,0),(w-1,1),\ldots,(0,w-1);
\]
all other incidence edges join two vertices already on this path.  Hence the
cut collar is connected.  Its Euler characteristic is that of a planar
surface by the computation \((*)\).  At even width and odd \(i\), the outer
face orbit identifies the two ends of the displayed path, producing exactly
the single relation among the \(p^2\) candidate disks; omitting the
lexicographically last disk breaks that relation.  Thus cutting the stated
disks removes every new one-handle and leaves a planar collar.  Inducting
over slabs shows that their boundaries are independent and form a complete
meridian system.  Their number is \(g(n,w)\), hence their span is
Lagrangian.

We construct the duals inductively in the longitudinal direction.  The
preceding cut-collar argument says more precisely that adjoining the terminal
slab removes \(2r_n\) disks from the old terminal cap region and joins their
boundary circles in \(r_n\) orientation-preserving pairs.  No other old face
orbit changes.  Gluing one such pair is an orientable one-handle attachment;
performing all pairs gives the relative homeomorphism
\[
 (\Sigma_{n+1,w},\text{old ribbon body})
 \cong
 (\Sigma_{n,w}\mathbin\# r_n(S^1\!\times S^1),
  \Sigma_{n,w}\setminus\text{terminal cap disks}),\qquad
 r_n=g(n+1,w)-g(n,w).
\]
This supplies the claimed relative connected-sum decomposition explicitly:
the meridian of each summand is the boundary of its displayed co-core disk,
and the path across its paired boundary circles is a dual longitude.
Cycles supported in the old ribbon body are
unchanged, their intersections are unchanged, and the old homology therefore
embeds as a symplectic subspace.  Its symplectic orthogonal complement is the
homology carried by the new slab collars.  The new classes in
\cref{eq:Ais} lie in that complement and form a Lagrangian there.

Retain the already chosen old pairs.  In the new complement choose preliminary
duals \(\widetilde B\) with
\(A_{\rm new}^T\Omega\widetilde B=I\).  Put
\(C=\widetilde B^T\Omega\widetilde B\), let \(U\) be the strictly upper
triangular matrix with the same upper triangle as \(C\), and set
\(B_{\rm new}=\widetilde B+A_{\rm new}U\).  Then
\[
 A_{\rm new}^T\Omega B_{\rm new}=I,
 \qquad B_{\rm new}^T\Omega B_{\rm new}=0,
\]
and orthogonality to the retained old pairs is automatic.  This proves by
induction that adding a terminal slice leaves every old pair fixed and only
appends new pairs.  The preliminary duals may be represented by geometric
longitudes in the corresponding slab collars; the triangular correction
adds only new co-core meridians.  Thus every new pair remains supported in
its slab collar, which gives filtration adaptation as well as nesting.
\end{proof}

If a geometric quadratic refinement is chosen as the affine origin, the
linear correction \(q_{\mathrm{geom}}(Se_j)\) on the canonical generators
must be transported with the basis.  The results below include this
correction; setting it silently to zero is not permitted.

\section{No spin-structure overhead on grid strips}

At a transverse separator, a partial even subgraph exposes the even frontier
mask
\begin{equation}\label{eq:Vw}
 V_w=\left\{m\in\Ftwo^{w^2}:\sum_{v\in\Lambda_w}m_v=0\right\},
 \qquad d_w=|V_w|=2^{w^2-1}.
\end{equation}

\begin{lemma}[grid relative exactness; \proofpending]\label{lem:grid-exact}
Move the transverse separator through the co-cores in the order of
\cref{lem:canonical}.  At every pair cut and every midpoint cut inside a
handle it meets exactly \(w^2\) labelled grid strands.  The cut satisfies
H1--H3 with trace space \(V_w\).
\end{lemma}

\begin{proof}
Moving through one checkerboard co-core replaces two intersection points on
one square by the other two and neither creates nor removes a longitudinal
strand.  Hence the labelled separator size stays \(w^2\).

Write \(G_L(\Gamma)\) and \(G_R(\Gamma)\) for the closed ribbon subgraphs on
the two sides of the moving separator \(\Gamma\), including its \(w^2\)
marked vertices, and put
\[
 P_\nu(\Gamma)=\{c\in C_1(G_\nu;\Ftwo):
                  \supp\partial c\subseteq\Gamma\},\qquad
 \tau_\nu c=(\partial c)|_\Gamma,\quad \nu=L,R.
\]
Choose the lexicographic spanning tree of the marked vertices in the collar
and let \(c_\nu(m)\) be its unique \(m\)-join.  These are the relative-chain
spaces, trace maps, and completions used below.

For a canonical generator \(u\), represent its character by the cellular
Poincar\'e-dual cocycle \(\gamma_u\): on a grid edge it is the mod-two
intersection of that edge band with the fixed meridian or longitude arc in
the slab collar.  If \(u\) is behind \(\Gamma\), every closed cycle in
\(G_R(\Gamma)\) has zero intersection with \(u\); hence
\(\gamma_u\) has zero period on \(\ker\tau_R\).  Choose a collar base vertex
\(v_*\), paths \(p_v\subset G_R\) from \(v_*\) to every right vertex, and set
\begin{equation}\label{eq:explicit-s-phi}
 s_u(v)=\langle\gamma_u,p_v\rangle,\qquad
 \phi_u(m)=\langle\gamma_u,c_R(m)\rangle+\langle s_u,m\rangle .
\end{equation}
Zero periods make \(s_u\) path independent.  For \(c\in P_R\), the chain
\(c+c_R(\tau_Rc)\) has zero trace, so
\[
 \langle\gamma_u,c\rangle
 =\langle\delta s_u,c\rangle+\phi_u(\tau_Rc).
\]
This is H1 with explicit \(s_u,\phi_u\).  A generator ahead of the separator
uses the identical construction in \(G_L\).

For H2, put
\[
 u=h_L(A_L,m),\quad v=h_R(A_R,m),\quad
 z(m)=\pi\bigl(c_L(m)+c_R(m)\bigr).
\]
At a pair cut, the slabwise connected-sum decomposition makes the left and
right handle subspaces symplectically orthogonal.  With \(B\) denoting the
intersection pairing, polarization gives the explicit split
\begin{align}
 q_0(u+v+z(m))
 &=\underbrace{q_0(u)+B(u,z(m))}_{Q_L(A_L,m)}
  +\underbrace{q_0(v)+B(v,z(m))}_{Q_R(A_R,m)}
  +\underbrace{q_0(z(m))}_{\kappa(m)}.              \label{eq:grid-H2}
\end{align}
Linear affine-origin corrections use \eqref{eq:explicit-s-phi} and therefore
alter only the two side terms and a mask function.

At the midpoint of handle \(i\), take \(a_i\) to be its co-core meridian and
choose the preliminary \(b_i\) as the longitude drawn in
\cref{fig:slab-cocores}.  On the exposed half, that longitude cocycle is
\(\delta s_i\), where \(s_i\) is the indicator of the vertices on the far
side of the co-core.  Thus
\begin{equation}\label{eq:stokes}
 \rho_i(m)=\langle\delta s_i,c\rangle
          =\langle s_i,\partial c\rangle
          =\sum_{v\in\Gamma}s_i(v)m_v .
\end{equation}
This is independent of the exposed chain \(c\) and proves H3.  In
\eqref{eq:grid-H2}, the only new cross-polarization at this midpoint is the
current \(a_i b_i\) term; its un-emitted value is \(\rho_i(m)\), so the term
is assigned to the \(a_i\) factor without enlarging the trace.

Finally consider the triangular correction in \cref{lem:canonical}.
Replacing \(\widetilde b_j\) by
\(b_j=\widetilde b_j+\sum_kU_{kj}a_k\) adds the corresponding cocycles.
The functions in \eqref{eq:explicit-s-phi} add in \(\Ftwo\), and all
\(a_k\) in a correction belong to the same newly attached slab complement.
Thus no corrected cocycle acquires a period on the opposite relative-cycle
space.  The pairing \(B\), and hence \eqref{eq:grid-H2}, is invariant under
the symplectic correction.  H1--H3 therefore survive the passage from the
geometric preliminary longitudes to the canonical basis.
\end{proof}

\begin{theorem}[all-sector grid rank bound; \proofpending]\label{thm:grid-upper}
For every \(n,w\geq2\), the complete pre-Arf tensor of the checkerboard
embedding of \(G_{n,w}\), in the canonical binary ordering, has every
flattening rank over \(\operatorname{Frac}\mathbb Z[t_e:e\in E]\) at most
\[
 d_w=2^{w^2-1}.
\]
The bound holds both between handle pairs and inside every handle.  It is
valid for arbitrary nonuniform edge weights and before the Arf sum.
Consequently a tensor train of this bond exists over the fraction field.
\end{theorem}

\begin{proof}
Apply \cref{thm:separator} and \cref{lem:grid-exact} at every successive
separator.  The final assertion follows from the characterization of
tensor-train ranks by successive flattening ranks.
\end{proof}

The following is the candidate denominator-free refinement and is not used
in the preceding theorem.  Let
\(\mathcal S_0,\ldots,\mathcal S_{2g}\) be the successive pair and half-handle
separators, with the mask labels identified with \(V_w\), and let \(E_j\) be
the edge block between \(\mathcal S_{j-1}\) and \(\mathcal S_j\).  For the local
spin bit \(\varepsilon_j\), define
\begin{equation}\label{eq:cores}
 A_j(\varepsilon_j)[m,m']
 =\sum_{\substack{S\subseteq E_j\\
       \partial_-S=m,\ \partial_+S=m'}}
 (-1)^{Q_j(S,m,m')+\varepsilon_jH_j(S,m,m')}
 \prod_{e\in S}t_e .
\end{equation}
The functions \(Q_j,H_j\) must be obtained from one global cochain gauge,
one compatible system of completion trees, and a phase assignment that
telescopes under gluing.  Defining them independently from the two-cut
factorizations is insufficient.  Proving that the edge blocks partition the
graph and that multiplication reproduces every term of \cref{eq:F} exactly
once is the remaining polynomial-ring obligation.  Conditional on that
obligation,
\begin{equation}\label{eq:allq-mps}
 \cF_{n,w}(\lambda)=
 \ell^TA_1(\lambda_{a_1})A_2(\lambda_{b_1})\cdots
 A_{2g}(\lambda_{b_g})r.
\end{equation}
the representation supports evaluation of one sector, Arf-weighted
contraction, partial sector sums, Walsh projection, and fixing selected
twists.  It does not permit all \(4^g\) values to be explicitly output in
polynomial time.

There are \(O(nw^2)\) binary coordinates.  A direct dense representation of
such denominator-free cores would take \(O(nw^2d_w^2)\) ring entries and a fixed boundary contraction
takes the same order of dense ring operations.  For fixed \(w\), both are
linear in \(n\).  These are separator-complexity bounds, not polynomial bounds
when \(w\) grows with \(n\).

\section{Generic minimality of the carrier}

With \(R_\infty(w)\) defined in \cref{eq:Rinf},
\cref{thm:grid-upper} targets \(R_\infty(w)\leq d_w\).  We next formulate
the matching generic lower bound and isolate its remaining proof obligations.
The proposed proof uses two different five-layer
encoders separated by a two-slab diagonal buffer; attempting to reflect one
encoder is false and is recorded in \cref{sec:failure}.

\subsection{One-sided encoders}

Let \(\mathcal G_w=G_{5,w}\) with terminal slice
\(S_w=\{(4,y,z):0\leq y,z<w\}\).  Let \(e_x(x,y,z)\), \(e_y(x,y,z)\), and
\(e_z(x,y,z)\) denote the positive coordinate edges.  Define the gauge tree
\begin{equation}\label{eq:T0}
 T_w^0=\{\text{all }e_x\}
 \cup\{e_y(0,y,z):0\leq y<w-1\}
 \cup\{e_z(0,0,z):0\leq z<w-1\}.
\end{equation}
Gauging homology cochains to vanish on \(T_w^0\), every chord is labelled by
the homology class of its fundamental cycle.

The explicit edge recurrences for the normal and opposite-phase encoder
trees are given in \cref{app:encoders}.  Their relevant property is the
following lemma.

\begin{lemma}[two encoder lemma; \proofpending]\label{lem:encoders}
For every \(w\geq4\), there are spanning trees \(T_w^+\) in the normal
five-layer checkerboard phase and \(T_w^-\) in the opposite five-layer phase
such that the terminal maps
\begin{equation}\label{eq:terminal-map}
 V(S_w)\longrightarrow H_1(\Sigma_{5,w};\Ftwo),\qquad
 m\longmapsto[\text{the unique tree }m\text{-join, completed at }S_w]
\end{equation}
are injective.  The construction is recursive in \(w\) and retains exactly
\(w^2-1\) independent homology chords.
\end{lemma}

\begin{proof}
The edge recurrences are in \cref{app:encoders}.  The four symbolic shell
inductions---normal and opposite, each at odd and even new width---are proved
in \cref{app:encoder-induction}.  They establish two invariants for every
width, rather than extrapolating from a finite check.  First, deleting the
\(w^2-1\) retained chords leaves a forest with exactly one terminal in every
component, so its terminal-cut columns form a basis of \(V(S_w)\).  Second,
after deletion of the duals of the gauge tree and retained chords, the
face-dual remainder is connected.  The tree--cotree criterion then makes the
corresponding fundamental-cycle homology labels independent.

The same appendix treats the exceptional normal set and the single
exceptional opposite chord at even width.  In the latter case the unique
dual-cut relation and terminal-cut column have disjoint support, so the
matrix determinant lemma preserves determinant one over \(\Ftwo\).  Hence
the terminal basis maps injectively to homology in both checkerboard phases,
which is \cref{eq:terminal-map}.
\end{proof}

\subsection{Buffered factor ranks}

Work in \(G_{11,w}\) and cut after the canonical handles in the first four
slabs.  At this cut, \cref{eq:XY} has the form
\begin{equation}\label{eq:XDY}
 \cF(\lambda_L,\lambda_R)=X(\lambda_L,m)D(m)Y(m,\lambda_R),
 \qquad m\in V_w,
\end{equation}
where \(D\) is a nonzero diagonal sign matrix.

\begin{lemma}[reachability and observability; \proofpending]\label{lem:XYfull}
Over the fraction field of the independent edge variables, \(X\) has column
rank \(d_w\) and \(Y\) has row rank \(d_w\).
\end{lemma}

\begin{proof}
On layers \(0,\ldots,4\), retain \(T_w^+\) and set all other left edge
variables to zero.  A tree has a unique partial even subgraph for each
terminal mask.  By \cref{lem:encoders}, the completed homology labels are
distinct.  Walsh transforming the left spin coordinates produces a square
character table of these labels multiplied by nonzero monomial and sign
diagonals.  Hence \(X\) has full column rank.

The last five layers \(6,\ldots,10\), after longitudinal reflection, have the
opposite checkerboard phase.  Retain \(T_w^-\).  Its terminal injectivity and
a Walsh transform give a full-rank suffix factor at slice \(6\).

Between slices \(4\) and \(6\), set all transverse edges to zero and retain
the two longitudinal edges at every \((y,z)\).  For each even mask \(m\), the
unique buffer configuration uses both longitudinal edges precisely at the
marked vertices.  The buffer transfer is therefore diagonal on \(V_w\), with
nonzero monomial diagonal.  It preserves the suffix rank, so \(Y\) has full
row rank.
\end{proof}

\begin{theorem}[generic minimality; \proofpending]\label{thm:grid-tight}
For every \(w\geq3\) and generic independent nonuniform edge weights,
\[
 R_\infty(w)=2^{w^2-1}.
\]
For every \(w\geq4\), the equality is already attained at a canonical pair
cut of \(G_{11,w}\).  Width three has a separate explicit \(256\times256\)
paired-cycle certificate.
\end{theorem}

\begin{proof}
For \(w\geq4\), choose a left inverse \(L\) of \(X\) and a right inverse
\(R\) of \(Y\) over the fraction field.  By \cref{eq:XDY},
\(L\cF R=D\), so the selected flattening has rank at least \(d_w\).  The
upper bound in \cref{thm:grid-upper} gives equality.

Every matrix entry is polynomial in the independent edge variables.  Thus a
full-rank minor over the fraction field is a nonzero real polynomial, whose
nonvanishing locus is Zariski open.  A nonzero real polynomial cannot vanish
on the open cube \(0<t_e<1\), so the same rank holds on a nonempty open set
of physical ferromagnetic high-temperature weights.

For \(w=3\), use \(G_{10,3}\) and cut after canonical handle five.  Order an
unoriented edge by first sorting its endpoints lexicographically and then
sorting the resulting endpoint pairs.  Retain the spanning-tree indices
\[
\begin{split}
T=\{&0,4,6,7,9,10,13,16,20,21,23,25,26,27,29,30,36,37,38,39,41,45,\\
&48,49,50,51,57,59,60,62,65,67,68,70,73,78,79,80,83,86,88,89,91,\\
&94,95,99,100,101,104,107,115,116,121,122,125,128,131,133,135,137,\\
&138,139,143,146,148,149,151,154,155,161,164,165,170,173,175,176,\\
&181,182,185,187,188,189,192,193,194,195,196,198,200\}
\end{split}
\]
and the chord indices
\[
 C=(44,52,103,110,112,118,120,162).
\]
All other edge weights are zero and retained weights are one.  The set
\(T\) has \(90-1=89\) edges and is connected; adjoining \(C\) therefore
gives cycle dimension eight.  In the ordered fundamental-cycle basis
specified by \(C\), select left canonical coordinates \(2,\ldots,9\) and all
eight right coordinates.  The two projection matrices are
\[
L=\begin{pmatrix}
1&0&0&0&1&0&0&0\\1&0&0&1&1&1&1&0\\1&0&0&0&0&0&0&0\\
0&0&0&1&0&1&1&1\\0&0&0&0&0&1&1&1\\1&1&0&0&0&0&0&0\\
1&1&0&1&0&1&0&0\\0&0&1&0&0&1&1&1
\end{pmatrix},\quad
R=\begin{pmatrix}
0&0&1&0&0&1&1&1\\1&1&0&1&0&1&0&0\\1&1&0&1&0&1&1&0\\
0&0&0&0&0&1&1&1\\0&0&0&1&0&1&1&1\\1&0&0&1&0&1&1&0\\
1&0&0&1&1&1&1&1\\0&0&0&0&0&0&0&1
\end{pmatrix}.
\]
Row reduction over \(\Ftwo\) gives
\(\det L=\det R=1\).  Hence both homology projections are isomorphisms.
Character orthogonality, applied on the two sides, gives the primary exact
factorization
\[
 W_LMW_R^T
 =2^{16}P_L\operatorname{diag}_{z\in\Ftwo^8}
       \bigl((-1)^{q_0(z)}\bigr)P_R .
\]
Both Walsh matrices and both permutations are invertible, so
\(\rank M=2^8=256\).  In the unnormalised original coordinates the
determinant has magnitude \(2^{8\cdot256}=2^{2048}\).  The modular values
\(409643880\pmod {1000000007}\) and
\(276933753\pmod {1000000009}\) are independent audits of this symbolic
argument, not its basis.

Under
the nested basis of \cref{lem:canonical}, append an eleventh slice, specialize
all newly incident edge variables to zero, and fix the newly appended sector
bits.  The same \(G_{10,3}\) minor is then a subminor of the \(G_{11,3}\)
flattening: every surviving chain lies in the old ribbon body, whereas every
new canonical pair is supported in the appended collar.  Thus \(n=11\) is
also sufficient at width three, which proves the uniform-length statement.
\end{proof}

The theorem does not imply that a chosen minor remains nonzero on the
homogeneous anisotropic locus \((t_x,t_y,t_z)\), on the isotropic line, at
every physical temperature, or at a critical point.  These are separate
algebraic restrictions and remain open here.

\section{A non-grid sharpness family}

The H3 condition in \cref{thm:separator} is genuinely stronger than pair-cut
factorization.  Take the toroidal cellular rotation of \(K_{3,3}\), with
bipartition \(0,1,2\mid3,4,5\), and rotation
\[
 \sigma_0=\sigma_1=\sigma_2=(3,4,5),\qquad
 \sigma_3=\sigma_4=\sigma_5=(0,1,2).
\]
The face permutation has the three hexagonal orbits
obtained from \(\varphi=\sigma\alpha\); hence \(V-E+F=6-9+3=0\) and the
surface is a torus.  Choose left and right port edges \((0,3)\) and \((1,5)\).
For an oriented edge-two-sum, delete the right port in copy \(j\) and the
left port in copy \(j+1\), identify \(1_j\sim0_{j+1}\) and
\(5_j\sim3_{j+1}\), and concatenate the two cyclic orders beginning
immediately after the deleted darts, reversing the order at the second
endpoint.  This is the boundary-orientation convention for gluing the two
port disks.  Iterating it defines \(C_r\).

\begin{theorem}[sharp internal obstruction; \proofpending]\label{thm:k33}
The surface of \(C_r\) has genus \(r\).  For arbitrary nonuniform weights,
the complete pre-Arf tensor has four-state handle-site bond at most two and
binary-coordinate bond at most four.  For generic independent weights every
nontrivial pair cut has rank exactly two, while every internal cut of a
non-end handle has rank exactly four.  The latter rank remains four under all
24 local affine symplectic relabellings of that handle.
\end{theorem}

\begin{proof}
Deleting a port ribbon leaves one disk in each torus; the stated
orientation-reversing boundary identification is their connected sum.
Induction gives genus \(r\).  Between
gadgets the separator has two marked points, with attainable masks \(00\) and
\(11\).  Handles behind and ahead have disjoint support, so H1--H2 hold and
\cref{thm:separator} gives pair bond two.  Splitting each handle gives the
unconditional binary bond four.

We print exact lower-bound witnesses.  Order the surviving unoriented edges
lexicographically by their ordered endpoints, work in
\(\mathbb F_p\), \(p=1000000007\), and set
\[
 t_e=2+104729(e+1)\pmod p.
\]
For two gadgets set the two external port weights (edge indices \(0,13\))
to zero.  At the pair cut, rows and columns \(0,1\) give
\[
 M_2=\begin{pmatrix}
580027348&972921000\\846017352&721285802
\end{pmatrix},\qquad
\det M_2=592118981\ne0\pmod p.
\]
For three gadgets set external port indices \(0,20\) to zero.  At the
internal cut after binary coordinate three, rows \(0,1,4,5\) and columns
\(0,1,2,3\) give
\[
M_3=\begin{pmatrix}
162853324&832578617&826262699&661314214\\
899577535&497270790&261851776&256133602\\
218465529&640022455&220866473&792750087\\
850525458&627258512&600657840&119903730
\end{pmatrix},
\qquad\det M_3=839109772\ne0\pmod p.
\]
These displayed matrices are independently checkable finite-field
specializations of integer-polynomial minors, so nonvanishing proves that
the symbolic minors are not identically zero.

At an arbitrary pair cut, zero every edge outside the two adjacent gadgets
and fix all exterior spin coordinates to \(00\); the surviving even
subgraphs and their signs are literally those defining \(M_2\).  At an
internal cut do the same outside the centered three-gadget block, also
zeroing its standalone outer ports; the surviving submatrix is \(M_3\).
Thus the minors embed unchanged for every chain length.

Finally, an affine symplectic relabelling in genus one is an arbitrary
permutation of the four quadratic refinements.  The determinants of the
same selected rank-four witness, after rank-revealing row selection, are:
\[
\begin{array}{c|r@{\qquad}c|r@{\qquad}c|r}
0123&839109772&0132&818787838&0213&192657990\\
0231&369943708&0312&465239613&0321&723491933\\
1023&818787838&1032&839109772&1203&465239613\\
1230&723491933&1302&192657990&1320&369943708\\
2013&369943708&2031&192657990&2103&723491933\\
2130&465239613&2301&839109772&2310&818787838\\
3012&723491933&3021&465239613&3102&369943708\\
3120&192657990&3201&818787838&3210&839109772
\end{array}\pmod p .
\]
Every entry is nonzero.  Hence no local affine symplectic relabelling
restores the H3 conclusion.
\end{proof}

This family separates two phenomena: separator compression can be exact at
handle-pair cuts while a binary ordering needs an extra local trace bit.  It
also prevents the grid midpoint argument from being mistaken for a formal
consequence of pair-cut compression alone.

\section{An all-sector algorithmic application}

Write the four-state version of \cref{eq:allq-mps} as
\[
 \cF(s_1,\ldots,s_g)=\ell^TA_1(s_1)\cdots A_g(s_g)r,
 \qquad s_i\in\Ftwo^2,
\]
and let \(\omega_i:\Ftwo^2\to R\) be arbitrary product-form sector weights.
For each handle \(i\) and Walsh character \(c\in\Ftwo^2\), define
\begin{equation}\label{eq:marginal}
 M_i(c)=\sum_{s\in(\Ftwo^2)^g}
   (-1)^{c\cdot s_i}\cF(s_1,\ldots,s_g)
   \prod_{j\ne i}\omega_j(s_j).
\end{equation}

\begin{theorem}[all single-handle Walsh marginals; \proved]
\label{thm:marginals}
Given an all-sector tensor train of bond at most \(d\), all \(4g\) quantities
in \cref{eq:marginal} can be computed with \(O(4gd^2)\) dense ring operations
and \(O(gd)\) environment storage, in addition to the cores.
\end{theorem}

\begin{proof}
Put
\(B_i=\sum_s\omega_i(s)A_i(s)\), and compute environments
\[
 L_0=\ell^T,\quad L_{i+1}=L_iB_i,
 \qquad
 R_g=r,\quad R_i=B_iR_{i+1}.
\]
Then
\begin{equation}\label{eq:marginal-contract}
 M_i(c)=\sum_{s\in\Ftwo^2}(-1)^{c\cdot s}
        L_iA_i(s)R_{i+1}.
\end{equation}
Expanding the products selects one local state at every handle and gives
exactly \cref{eq:marginal}.  The two sweeps and all four local contractions
per handle have the stated cost.
\end{proof}

An explicit sector-list calculation must inspect \(4^g\) values before
applying arbitrary product weights.  The comparison is for the \(4g\)
all-family observables \cref{eq:marginal}, not for computing only the physical
partition function \cref{eq:arf}; ordinary transfer already has the same
separator dependence for that single contraction.

\section{Embedding and coordinate robustness}

\begin{theorem}[robust class; \proved]\label{thm:robust}
The separator bound is invariant under filtration-preserving surface
homeomorphisms; coboundary changes of character representatives; relative-
boundary changes of completion chains; explicitly transported affine changes
of quadratic origin; and invertible coordinate changes supported entirely on
one side of each pair cut.  At an internal cut, the coordinate change must
also preserve H3.  The bound is not invariant under an arbitrary rotation
system or an arbitrary symplectic mixing across a declared cut.
\end{theorem}

\begin{proof}
A filtration-preserving homeomorphism relabels partial chains and masks.
Coboundary changes are mask-diagonal by \cref{eq:stokes}.  Changing a
completion alters the phase by left-only, right-only, and mask-only terms.
Affine corrections are H1 characters.  A coordinate change confined to one
side multiplies the flattening by an invertible matrix on that side.  These
operations preserve \cref{eq:XY}.  A transformation mixing an internal cut
need not preserve the trace property, as \cref{thm:k33} shows.
\end{proof}

\proved: Stabilizing the surface in a disk disjoint from the graph makes the pre-Arf
tensor independent of the two new bits; the new four-state site is the
rank-one vector \((1,1,1,1)\).  Such a stabilization is noncellular.  Indeed,
for a connected cellular embedding the cellular chain complex makes
\(H_1(G;\Ftwo)\to H_1(\Sigma;\Ftwo)\) surjective, so no surface handle is
homologically unused by the graph.

\proved: Arbitrary cellular rotations are different.  The cube graph
\(P_2\square P_2\square P_2\) has orientable rotation systems of genera zero,
one, and two.  Exhausting the \(2^8\) choices of local orientation relative to
fixed cyclic orders gives respectively \(2,54,200\) rotations.  The genus-zero
and genus-two pre-Arf families have sizes one and sixteen and cannot be
related by an invertible map in one fixed sector space.  Their normalized Arf
sums nevertheless give the same embedding-independent even-subgraph
polynomial, as required by \cref{eq:arf}.

\section{Exact validation and failure ledger}

The theorems above are algebraic; finite computation is not used as a
substitute for the arbitrary-width proofs.  It is used as an independent
coordinate and implementation firewall.

\subsection{Validation}

\certified: Over both primes \(1{,}000{,}000{,}007\) and
\(1{,}000{,}000{,}009\), exact scripts verify the following:

\begin{itemize}
\item \(S^T\Omega S=J\) over \(\Ftwo\), quadratic polarization, affine
  corrections, and basis-independent Arf signs for \(2\leq n,w\leq6\), with
  additional embedding stress tests at \(w=7,8\);
\item agreement of direct spin, even-subgraph, domain-wall, Kac--Ward, and
  compressed-transfer evaluations on the small graphs in the proof archive;
\item pair and internal ranks independently, including the canonical
  rank profiles in \cref{tab:ranks};
\item all \(4g\) Walsh marginals from \cref{eq:marginal-contract} against
  literal enumeration of \(4^g\) sectors for \(G_{6,3}\) and \(G_{7,3}\);
\item topology and nonzero minors for the \(K_{3,3}\) sharpness family, and
  the exhaustive 256-rotation census of the cube graph.
\end{itemize}

\begin{table}[ht]
\centering
\caption{Exact canonical binary rank profiles, identical over the two
primes.  The nonuniform weights are independently specialized; the other
columns are additional checks, not arbitrary-width homogeneous theorems.}
\label{tab:ranks}
\small
\begin{tabular}{@{}lll@{}}
\toprule
graph & weights & binary rank profile\\
\midrule
\(G_{10,3}\) & nonuniform &
\(2,4,8,16,32,64,128,256,256,256,128,64,32,16,8,4,2\)\\
\(G_{10,3}\) & anisotropic \((2,3,5)\) & same\\
\(G_{10,3}\) & isotropic \(t=2\) & same\\
\(G_{4,4}\) & nonuniform & \(2,4,8,16,32,16,8,4,2\)\\
\(G_{4,4}\) & anisotropic \((2,3,5)\) & same\\
\(G_{4,4}\) & isotropic \(t=2\) & same\\
\bottomrule
\end{tabular}
\end{table}

For the marginal test, \(G_{6,3}\) has genus five, \(1024\) sectors, and
pair ranks \(4,16,16,4\); \(G_{7,3}\) has genus six, \(4096\) sectors, and
pair ranks \(4,16,64,16,4\).  The exact tensor-train outputs equal literal
sector summation over both primes.

\subsection{Contained failures}\label{sec:failure}

Two failures constrain the argument: a noncanonical coordinate system
creates a spurious width-three internal rank \(512\), and reflection of the
normal encoder gives a rank-deficient right map.  Their full data and the
discarded inferences are moved to
\path{failure-ledger-supplement.tex}.  The present proof uses only canonical
coordinates and an independently defined opposite-phase encoder.

\section{Discussion}

The central structural target is:
\[
 \boxed{\text{complete spin-structure bond}
        =\text{ordinary even-frontier bond generically}.}
\]
If the pending grid and encoder lemmas are completed, the upper equality
means that repeated handle attachment, Arf/Walsh phases,
and cochain gauge changes introduce no additional exact memory once the
frontier mask is retained.  The generic lower equality means that no
universal exact compression of this complete tensor below the transverse
carrier can exist for arbitrary independent weights.

This result refines, rather than bypasses, the usual width barrier.  At fixed
\(w\) it is a rigorous strip compression, and for \(w=L\) its exact bond is
area exponential.  It does not conflict with complexity results for richer
families of nonplanar Ising instances, which allow selectable couplings and
different input classes \cite{Istrail2000}.  Nor does it settle the
homogeneous simple-cubic sequence: an independent-weight minor can vanish
identically after a homogeneous specialization.

Three concrete questions remain.  First, does generic minimality persist on
the homogeneous anisotropic or isotropic loci?  Second, can the abstract
separator theorem be applied to other surface decompositions with trace
dimension smaller than the ordinary spin frontier?  Third, can a different
representation---not the complete pre-Arf family studied here---replace the
area carrier by a structure stable under the cubic thermodynamic limit?
None is answered by the present work.

\appendix

\section{Explicit recursive encoder data}\label{app:encoders}

This appendix freezes the constructive content of \cref{lem:encoders}.  An
edge \(e_\alpha(x,y,z)\) starts at \((x,y,z)\) and points in coordinate
\(\alpha\in\{x,y,z\}\).  Empty ranges contribute no edges.  The width-four
normal tree is the 79-edge base stored verbatim in the ancillary certificate;
equivalently it is the edge-index set
\begin{verbatim}
1,2,5,7,8,9,10,13,15,16,19,21,23,24,25,26,27,29,30,32,
34,36,37,39,42,44,45,48,50,51,53,56,57,59,61,64,67,69,70,72,
74,76,78,79,82,83,85,88,90,92,93,96,99,100,101,104,107,110,
112,114,116,118,119,122,125,130,139,147,152,154,158,159,162,
163,167,168,171,172,177
\end{verbatim}
in the repository convention that sorts grid edges lexicographically by
their ordered endpoints.  The verifier checks that this list has 79 entries;
the display is data, not a claim inferred from the printed line wrapping.

For the normal extension from old width \(o\) to \(o+1\), first include all
new \(x\)-rails
\[
 e_x(r,o,z),\quad 0\leq z\leq o,
 \qquad
 e_x(r,y,o),\quad0\leq y<o,
 \qquad r=0,1,2.
\]
If \(o\) is even, add
\begin{align*}
&e_y(0,o-1,z)\ (0\leq z\leq o),\quad e_y(0,0,o),\quad
 e_y(0,y,o)\ (y\text{ odd}),\\
&e_z(0,y,o-1)\ (2\leq y<o,\ y\text{ even}),\quad e_y(1,o-2,o),\\
&e_x(3,o,o),\quad e_x(3,y,o)\ (2\leq y<o,\ y\text{ even}),\\
&e_y(4,o-1,z),e_y(4,z,o),e_z(4,o,z)\ (0\leq z<o,\ z\text{ even}),\\
&e_z(4,0,o-1).
\end{align*}
If \(o\) is odd, add
\begin{align*}
&e_y(0,o-1,z)\ (3\leq z\leq o,\ z\text{ odd}),\\
&e_y(0,y,o)\ (y=0,o-2\text{ or }2\leq y<o\text{ even}),\\
&e_z(0,o,0),\quad e_z(0,o,z)\ (z\text{ odd}),\quad e_y(1,o-1,1),\\
&e_z(1,y,o-1)\ (y\text{ odd}),\quad
 e_x(3,o,z)\ (z\text{ even or }z=o),\quad e_x(3,0,o),\\
&e_y(4,o-1,z),e_y(4,z,o)\ (z\text{ odd}),\quad
 e_z(4,y,o-1)\ (y\text{ odd}).
\end{align*}
Each extension contains \(10o+5\) edges, one per new vertex.

For the normal phase put \(k=\lfloor w/2\rfloor\) and define the exceptional
tree chords
\begin{align}\label{eq:Xplus}
 X_w^+={}&\{e_z(0,3,2)\}
 \cup\{e_z(0,y,0):y=5,7,\ldots,2k-1\}\notag\\
 &\cup\{e_z(1,2a+1,2b):1\leq a\leq k-3,\ a+2\leq b\leq k-1\},
\end{align}
and retain
\begin{equation}\label{eq:Pplus}
 P_w^+=(T_w^+\setminus T_w^0)\setminus X_w^+.
\end{equation}
Empty ranges are again omitted.  Direct subtraction from the preceding
recurrence gives \(|P_w^+|=w^2-1\).

The opposite width-four base is obtained from the normal base by the 13-edge
exchange in \cref{tab:exchange}.  For all later extensions, first include the
new \(x\)-rails at \(x=1,2\).  If \(o\) is odd, add
\begin{align*}
&e_x(0,y,o)\ (y\ne1),\quad e_x(0,o,z),\quad
 e_x(3,y,o)\ (y=0\text{ or }y\text{ odd}),\quad e_x(3,o,0),e_x(3,o,o),\\
&e_y(0,y,o)\ (y\text{ even},y<o-1),\quad e_y(0,o-1,z)\ (2\leq z\leq o),\\
&e_y(1,1,o),e_y(1,o-2,o),\quad e_y(4,y,o)\ (y\text{ odd}),
 e_y(4,o-1,z)\ (z\text{ odd}),\\
&e_z(0,y,o-1)\ (y\text{ odd}),\quad e_z(0,o,0),e_z(1,o,1),
 e_z(4,o,z)\ (z\text{ odd}).
\end{align*}
If \(o\geq6\) is even, add
\begin{align*}
&e_x(0,y,o),e_x(0,o,z),\quad e_x(3,o,z)\ (z\text{ odd or }z=o),\\
&e_y(0,y,o)\ (y<o-1),\quad
 e_y(0,o-1,z)\ (z=0,1,2\text{ or positive even}),\\
&e_y(4,y,o)\ (y\text{ even}),\quad e_y(4,o-1,z)\ (z\text{ even}),\\
&e_z(0,o,z)\ (z=2,4,\ldots,o-2),\quad e_z(1,0,o-1),
 e_z(4,y,o-1)\ (y\text{ even}).
\end{align*}
The \(o=4\) shell uses \(e_x(3,4,z)\) for \(z=1,2,4\), the \(x=4\)
shell positions \(z=0,3\), and \(e_z(0,4,0)\), exactly as frozen by the
ancillary verifier.

For the opposite phase define
\begin{equation}\label{eq:Xminus}
 X_w^-=\begin{cases}
 \varnothing,&w\text{ odd},\\
 \{e_z(0,3,2)\},&w=4,\\
 \{e_z(0,w-1,0)\},&w\geq6\text{ even},
 \end{cases}
 \qquad
 P_w^-=(T_w^-\setminus T_w^0)\setminus X_w^-.
\end{equation}
The recurrences give \(|P_w^-|=w^2-1\).  These are the retained chord sets
used in every subsequent primal and dual deletion; no implicit fitted subset
is being chosen.

\begin{table}[p]
\centering
\caption{Thirteen-edge exchange from the normal to the opposite width-four
base.  Each tuple denotes an unoriented edge.}
\label{tab:exchange}
\scriptsize
\begin{tabular}{@{}p{.47\textwidth}p{.47\textwidth}@{}}
\toprule
remove & add\\
\midrule
\((0,1,0)(1,1,0)\), \((0,1,3)(1,1,3)\),
\((0,2,1)(0,2,2)\), \((1,0,1)(1,1,1)\),
\((1,1,0)(1,1,1)\), \((1,1,2)(1,1,3)\),
\((1,2,2)(1,3,2)\), \((2,0,1)(2,0,2)\),
\((2,1,0)(2,2,0)\), \((2,1,3)(2,2,3)\),
\((3,2,1)(4,2,1)\), \((4,1,0)(4,1,1)\),
\((4,1,2)(4,1,3)\)
&
\((0,0,1)(0,1,1)\), \((0,1,0)(0,1,1)\),
\((0,1,2)(0,1,3)\), \((1,0,1)(1,0,2)\),
\((1,1,0)(1,2,0)\), \((1,1,3)(1,2,3)\),
\((1,2,1)(1,2,2)\), \((2,1,0)(2,1,1)\),
\((2,1,2)(2,1,3)\), \((2,2,2)(2,3,2)\),
\((3,1,0)(4,1,0)\), \((3,1,3)(4,1,3)\),
\((4,2,1)(4,2,2)\)\\
\bottomrule
\end{tabular}
\end{table}

\begin{table}[p]
\centering
\caption{Finite component types used in the parity induction of
\cref{lem:encoders}.  ``Old 0'' is an unchanged old terminal component.}
\label{tab:encoder-checks}
\scriptsize
\begin{tabular}{@{}p{.14\textwidth}p{.40\textwidth}p{.34\textwidth}@{}}
\toprule
case & terminal components after deleting retained chords
& breadth-first layer sizes of the contracted dual remainder\\
\midrule
normal \(w=4\) & 16 components, one per terminal & direct connected check\\
normal odd \(w\geq7\) & new: \(3(w-1)/2+1\) of size 1,
\((w-3)/2\) of size 9, one of size 13; old: \(w-3\) of size 4, one of
size 8, otherwise old 0 & \(1,3w-4,w-1,w-1,1\)\\
normal even \(w\geq6\) & new: \(3(w-2)/2\) of size 1,
\((w-2)/2\) of size 5, two of size 9, one of size 17; old:
\((w-4)/2\) of size 4, \((w-6)/2\) of size 8, otherwise old 0
& \(1,3w-6,w-2,w-2\)\\
opposite \(w=4\) & new sizes \(1,4,4,5,5,6,6\), one old size 4,
eight old 0 & direct connected check\\
opposite odd \(w\geq5\) & new: \(3(w-1)/2\) of size 1,
\((w-1)/2\) of size 5, one of size \(4w+1\); old: \((w-1)/2\) of size 4,
otherwise old 0 & \(1,2w-3,w-2,w-1\)\\
opposite even \(w\geq6\) & new: \(3(w-2)/2\) of size 1, one of size 4,
one of size 6, \(w/2\) of size 9; old: \(w-3\) of size 4, otherwise old 0
& \(1,2w-4,w-2,w-2,3\)\\
\bottomrule
\end{tabular}
\end{table}

\clearpage
\section{Symbolic shell induction for the encoder lemma}
\label{app:encoder-induction}

We give the symbolic shell induction intended to show that the recurrences
of \cref{app:encoders} preserve the two encoder invariants at arbitrary
width.  The finite-width programs described in the reproducibility section
are regression checks only.  The remaining base and exceptional-mode
obligations are stated at the end of the appendix.

Write \(o=W-1\) for the old and new widths.  The new transverse
L-shell has \(5(2o+1)=10W-5\) vertices.  Old components are contracted
before examining the shell.  Since each recurrence retains the old tree and
adds one parent edge for every new vertex, with no edge joining two already
reached vertices, the new graph is again a spanning tree.  All old faces
away from \(y=W-2\) and \(z=W-2\) keep their former dual component.  Thus a
shell can meet the old contracted dual remainder only along these two
boundary rows.  The edge lists in \cref{app:encoders} depend there only on
the parity of the running coordinate.  Consequently each induction step is
the repetition of the finite local parent patterns recorded below, plus the
unchanged old remainder.

\paragraph{The terminal-forest invariant.}
Delete the retained chords.  Following the rails from each new vertex until
the first terminal or contracted old component gives the component types in
the middle column of \cref{tab:encoder-checks}.  The following four sums show
that those types exhaust the new shell:
\begin{align*}
 W\text{ odd, normal: }&
 \left(\frac{3(W-1)}2+1\right)
 +9\frac{W-3}2+13+4(W-3)+8=10W-5,\\
 W\text{ even, normal: }&
 3\frac{W-2}2+5\frac{W-2}2+2\cdot9+17
 +4\frac{W-4}2+8\frac{W-6}2=10W-5,\\
 W\text{ odd, opposite: }&
 3\frac{W-1}2+5\frac{W-1}2+(4W+1)
 +4\frac{W-1}2=10W-5,\\
 W\text{ even, opposite: }&
 3\frac{W-2}2+4+6+9\frac W2+4(W-3)=10W-5.
\end{align*}
The summands occur in the order ``new components, modified old
components'' used in the table.  Their anchors are distinct: a new component
ends at its displayed shell terminal, while a modified old component retains
the unique terminal supplied by the induction hypothesis.  Counting the
unchanged old components gives exactly \(W^2\) anchors in every row.  The
four vertex identities therefore leave neither an anchor-free component nor
an unassigned shell vertex.  Because the graph is a forest, every component
contains exactly one terminal.  This proves the terminal invariant by
induction from the two width-four bases (with the exceptional normal
\(4\to5\) step checked directly from the displayed recurrence).

\paragraph{Normal dual parent map.}
Denote a boundary square by \((a;x,y,z)\), where \(a\) is its fixed
coordinate axis and \((x,y,z)\) its lower corner.  Contract separately each
connected old/bulk or changed/new face component after deleting
\((T_W^0\cup P_W)^*\).  For odd \(W\), the only nondirect shell arms are,
for even \(0\le j\le W-3\),
\[
 \begin{array}{ccccc}
 (2;3,j,W-2)&\longleftarrow&(0;0,j,W-2)
 &\longleftarrow&(2;1,j,W-2),\\
 (1;3,W-2,j)&\longleftarrow&(0;1,W-2,j)
 &\longleftarrow&(1;1,W-2,j).
 \end{array}
\]
The arrows point toward the parent.  The first arm with \(j=W-3\) has the
additional child \((0;0,W-3,W-2)\); every other contracted component meets
the outer component directly.  The two arm families have \((W-1)/2\)
members each.  Together with the unchanged direct bulk components they give
the breadth-first counts
\[
 1,\quad 3W-4,\quad W-1,\quad W-1,\quad1.
\]

For even \(W\), each odd \(1\le j\le W-3\) gives the parent tree
\[
 (1;0,W-2,j)\longleftarrow(0;1,j,W-2)
 \longleftarrow
 \{(2;0,j,W-2),(2;2,j,W-2)\}.
\]
The old component containing \((0;0,0,W-3)\) additionally parents the
\((W-2)/2\) leaves \((0;0,W-2,j)\) with odd \(j\); all other old components
meet the outer component.  Hence the layer counts are
\[
 1,\quad3W-6,\quad W-2,\quad W-2.
\]
In both parities, consecutive displayed squares share the unique transverse
edge left undeleted by the corresponding recurrence.  Splitting the two
boundary rows according to coordinate parity gives exactly the displayed
families; a corner belongs to the distinguished family, and an interior
boundary square belongs to one value of \(j\).  Every other cross-component
boundary edge is in \(T_W^0\cup P_W\), while every off-shell adjacency stays
inside an old component.  Thus every nonroot quotient component occurs once
and has one parent.  The quotient has one fewer displayed edge than vertices
and is a tree.  This proves connectedness of the normal dual remainder for
all \(W\).

\paragraph{Opposite dual parent map.}
The same contraction applied to the opposite recurrence admits the following
complete parent description.  This also corrects a layer-count transcription
error in the earlier working proof: the correct second-layer counts are
\(2W-3\) for odd \(W\) and \(2W-4\) for even \(W\).

Suppose first that \(W\geq5\) is odd, and put
\(r=W-2\).  The root is the component represented by
\(O=(0;0,0,r)\).  The components adjacent to \(O\) are
\begin{align*}
 A_j&=(0;0,0,j) &&(j=1,3,\ldots,r-2),\\
 B_j&=(0;0,1,j),&
 C_j&=(1;0,r,j),&
 D_j&=(1;2,r,j) &&(j=0,2,\ldots,r-1).
\end{align*}
Here and below every member of a displayed family has the displayed parent;
thus \(A_j,B_j,C_j,D_j\leftarrow O\).  The remaining parent edges are
\begin{align*}
 E_j=(0;0,r,j)&\leftarrow C_j &&(j=2,4,\ldots,r-1),\\
 F_j=(0;1,j,r)&\leftarrow B_{r-1} &&(j=0,2,\ldots,r-1),\\
 (2;0,j,r)&\leftarrow F_j,&
 (2;2,j,r)&\leftarrow F_j &&(j=0,2,\ldots,r-1).
\end{align*}
Consequently the breadth-first layers have cardinalities
\[
 1,\quad2W-3,\quad W-2,\quad W-1.
\]
The formula includes the base shell \(W=5\), where the layers are
\(1,7,3,4\).  In that base shell only, the first line is replaced by
\(E_0=(0;0,3,0)\leftarrow C_0\); its cardinality is unchanged.

Now suppose that \(W\geq6\) is even and again put \(r=W-2\).  The root is
\(O=(0;0,r,1)\).  Its children are
\begin{align*}
 A_j&=(0;0,0,j) &&(j=1,3,\ldots,r-1),\\
 B_j&=(0;0,1,j) &&(j=0,2,\ldots,r-2),\\
 C_j&=(1;3,r,j),&
 D_j&=(2;3,j,r) &&(j=1,3,\ldots,r-1).
\end{align*}
The next layer and its parents are
\begin{align*}
 E_j=(0;0,j,r)&\leftarrow D_j &&(j=3,5,\ldots,r-3),\\
 F_j=(0;1,r,j)&\leftarrow C_j &&(j=3,5,\ldots,r-1),\\
 G_1=(0;2,1,r)&\leftarrow D_1,&
 G_{r-1}=(0;2,r-1,r)&\leftarrow D_{r-1},\\
 H=(0;2,r,1)&\leftarrow C_1.&
\end{align*}
Empty ranges are omitted.  The third layer is
\begin{align*}
 I_1=(1;1,r,1)&\leftarrow H,&
 I_j=(1;1,r,j)&\leftarrow F_j &&(j=3,5,\ldots,r-1),\\
 J_1=(2;1,1,r)&\leftarrow G_1,&
 J_{r-1}=(2;1,r-1,r)&\leftarrow G_{r-1},\\
 J_j=(2;1,j,r)&\leftarrow E_j &&(j=3,5,\ldots,r-3),
\end{align*}
and the final three leaves are
\[
 (0;0,1,r)\leftarrow J_1,qquad
 (0;0,r-1,r)\leftarrow J_{r-1},qquad
 (0;1,r,1)\leftarrow I_1.
\]
The even-width layers therefore have sizes
\[
 1,\quad2W-4,\quad W-2,\quad W-2,\quad3.
\]

It remains to verify that these are the \emph{complete} quotient edge sets,
not merely spanning parent sets.  Two boundary squares in the list are
adjacent exactly when their descriptors share the transverse edge indicated
by the arrow.  Substitution in the odd and even edge recurrences of
\cref{app:encoders} shows that every arrow edge is outside
\(T_W^0\cup P_W^-\).  The other boundary adjacencies either join faces
already contracted into the same old/new component or use an edge of
\(T_W^0\cup P_W^-\).  Off-shell adjacencies remain inside an old component.
The displayed families are disjoint and exhaust the parity classes on both
boundary rows, so there are no omitted quotient vertices.  Hence the
quotient has exactly one displayed edge per nonroot component and is a tree.
This proves connectedness of the opposite dual remainder, conditional only
on the printed width-four base trace discussed below.

\paragraph{Exceptional modes.}
Let \(D_W\) be the face dual and let \(h_e\) denote the fundamental-cycle
homology label of a chord relative to \(T_W^0\).  For any chord set \(A\),
the cographic rank formula is
\begin{equation}\label{eq:tree-cotree-rank}
 \rank h(A)=|A|-c\bigl(D_W-(T_W^0\cup A)^*\bigr)+1.
\end{equation}

For the normal encoder, deletion with and without the exceptional set
\(X_W^+\) has the same dual-component partition.  Besides the large
component, its small components are
\begin{equation}\label{eq:normal-islands}
 \begin{split}
 I_3&=\{(0;0,y,2):0\le y\le2\},\\
 I_5&=\{(0;0,y,0):0\le y\le4\},\\
 I_{2,r}&=\{(0;0,2r-1,0),(0;0,2r,0)\},
       \qquad3\le r<\lfloor W/2\rfloor .
 \end{split}
\end{equation}
Indeed, along the \(x=0\) boundary row, the first exceptional tooth closes
\(I_3\), the first \(z=0\) tooth closes \(I_5\), and each subsequent
odd-row tooth closes the displayed pair.  Every other boundary edge of
these sets lies in \(T_W^0\cup X_W^+\), while the triangular \(x=1\) teeth
and all edges of \(P_W^+\) have both incident faces in the same displayed
set or in the large component.  Thus the component claim reduces to the
finite parity families in \cref{eq:normal-islands}; the still-unprinted
boundary-incidence table is the explicit check of this sentence.  Once that
check is supplied, applying
\cref{eq:tree-cotree-rank} to \(X_W^+\) and to
\(P_W^+\cup X_W^+\), and using connectedness for \(P_W^+\), shows
\begin{equation}\label{eq:normal-span-split}
 \operatorname{span}h(P_W^+)\cap\operatorname{span}h(X_W^+)=0.
\end{equation}

For the opposite encoder, \(X_W^-\) is empty at odd width.  At even
\(W\ge6\), put \(X=e_z(0,W-1,0)\).  Removing additionally \(X^*\) from the
connected dual remainder splits off the following face set:
\begin{align}\label{eq:opposite-cut}
 C_W={}&\{(1;x,y,0):x\in\{0,2\},\ 0\le y<W\}\notag\\
 &\cup\{(1;x,y,0):x\in\{1,3\},\ y\in\{0,W-1\}\}\notag\\
 &\cup\{(2;x,y,0):x\in\{0,2\},\ 0\le y\le W-2,\ y\text{ even}\}\notag\\
 &\cup\{(2;x,y,0):x\in\{1,3\},\ 0\le y\le W-2\}\notag\\
 &\cup\{(0;x,y,0):0\le x\le3,\ 1\le y\le W-3,\ y\text{ odd}\}\notag\\
 &\cup\{(0;4,y,0):0\le y\le W-2\}\notag\\
 &\cup\{(2;x,y,1):x\in\{1,3\},\ 1\le y\le W-3,\ y\text{ odd}\}.
\end{align}
Successive squares in each row of \cref{eq:opposite-cut} are joined by the
undeleted rails in the even recurrence, so \(C_W\) is connected.  Inspecting
its boundary shows that the non-gauge chord edges crossing the cut are
exactly
\[
 \{X\}\cup\{e_y(4,2j,1):0\le j<W/2\}.
\]
Gauge-tree edges have zero label.  The face-boundary relation for
\(C_W\) therefore gives the unique dependence
\begin{equation}\label{eq:opposite-h-relation}
 h_X=\sum_{j=0}^{W/2-1}h_{e_y(4,2j,1)}.
\end{equation}

It remains to compare terminal cuts.  Root the terminal slice at
\(s_0=(4,0,0)\), and for a tree edge \(e\) let \(u_e\in V(S_W)\) be the
indicator of the terminals on the component of \(T_W^-\setminus e\) not
containing \(s_0\).  Following the printed even-width rails shows that both
\(X\) and \(e_z(1,W-1,1)\) separate the sole terminal \((4,W-1,0)\).  Hence
\begin{equation}\label{eq:opposite-u-relation}
 u_X=u_{e_z(1,W-1,1)}.
\end{equation}
The edge on the right of \cref{eq:opposite-u-relation} does not occur in
\cref{eq:opposite-h-relation}.  At \(W=4\), the complete base trace below
gives instead
\begin{align*}
 h_X&=h_{e_y(4,0,1)}+h_{e_y(4,2,1)},\\
 u_X&=u_{e_z(1,2,1)}+u_{e_y(2,2,2)},
\end{align*}
again with disjoint supports.

Finally use the \(P_W^-\) terminal-cut columns and homology labels as bases.
Without \(X\), the terminal-to-homology matrix is the identity.  Adding
\(X\) is the rank-one update \(I+uv^T\), where \(u\) and \(v\) are the two
coordinate vectors just displayed.  Their supports are disjoint, so
\(v^Tu=0\) and
\[
 \det(I+uv^T)=1+v^Tu=1\quad\text{in }\Ftwo.
\]
Thus the exceptional edge does not reduce terminal injectivity.  The
opposite quotient parent map and both finite width-four base traces are
printed below.  The remaining manuscript-level gap in \cref{lem:encoders}
is the fully expanded boundary-incidence verification of
\cref{eq:normal-islands} and \cref{eq:opposite-cut}; until that bookkeeping is included,
the lemma remains marked \proofpending.

\section{Complete normal width-four encoder base trace}
\label{app:width4-normal-trace}
This appendix discharges the finite base case used by the normal shell induction.  Vertices and edge indices use the convention of \cref{app:encoders}.  A face $(a;x,y,z)$ is the unit square with fixed coordinate axis $a$ and lower corner $(x,y,z)$; $O$ denotes the nonsquare outer face.  Each row below is an explicit incidence check, so no program output is needed to infer connectivity.
The normal base contains 80 vertices and the 79 selected tree edges printed in \cref{app:encoders}.  Starting at $(0,0,0)$, the following table gives one parent edge for every other vertex.
\begin{longtable}{@{}rrr@{}}
\toprule child & parent & edge index\\ \midrule
\endfirsthead \toprule child & parent & edge index\\ \midrule \endhead
$(0,1,0)$ & $(0,0,0)$ & $1$\\
$(1,0,0)$ & $(0,0,0)$ & $2$\\
$(1,1,0)$ & $(0,1,0)$ & $13$\\
$(2,0,0)$ & $(1,0,0)$ & $42$\\
$(1,1,1)$ & $(1,1,0)$ & $51$\\
$(2,1,0)$ & $(1,1,0)$ & $53$\\
$(3,0,0)$ & $(2,0,0)$ & $82$\\
$(0,1,1)$ & $(1,1,1)$ & $16$\\
$(1,0,1)$ & $(1,1,1)$ & $44$\\
$(2,1,1)$ & $(1,1,1)$ & $56$\\
$(2,2,0)$ & $(2,1,0)$ & $92$\\
$(3,1,0)$ & $(2,1,0)$ & $93$\\
$(4,0,0)$ & $(3,0,0)$ & $122$\\
$(0,2,1)$ & $(0,1,1)$ & $15$\\
$(0,0,1)$ & $(1,0,1)$ & $5$\\
$(2,0,1)$ & $(1,0,1)$ & $45$\\
$(3,1,1)$ & $(2,1,1)$ & $96$\\
$(1,2,0)$ & $(2,2,0)$ & $64$\\
$(3,2,0)$ & $(2,2,0)$ & $104$\\
$(0,2,2)$ & $(0,2,1)$ & $25$\\
$(0,3,1)$ & $(0,2,1)$ & $26$\\
$(1,2,1)$ & $(0,2,1)$ & $27$\\
$(2,0,2)$ & $(2,0,1)$ & $83$\\
$(3,0,1)$ & $(2,0,1)$ & $85$\\
$(0,2,0)$ & $(1,2,0)$ & $24$\\
$(0,3,2)$ & $(0,2,2)$ & $29$\\
$(1,2,2)$ & $(0,2,2)$ & $30$\\
$(1,3,1)$ & $(0,3,1)$ & $36$\\
$(2,2,1)$ & $(1,2,1)$ & $67$\\
$(1,0,2)$ & $(2,0,2)$ & $48$\\
$(3,0,2)$ & $(2,0,2)$ & $88$\\
$(4,0,1)$ & $(3,0,1)$ & $125$\\
$(0,3,0)$ & $(0,2,0)$ & $23$\\
$(0,3,3)$ & $(0,3,2)$ & $37$\\
$(1,3,2)$ & $(1,2,2)$ & $69$\\
$(2,2,2)$ & $(1,2,2)$ & $70$\\
$(2,3,1)$ & $(1,3,1)$ & $76$\\
$(3,2,1)$ & $(2,2,1)$ & $107$\\
$(0,0,2)$ & $(1,0,2)$ & $8$\\
$(4,0,2)$ & $(4,0,1)$ & $162$\\
$(4,1,1)$ & $(4,0,1)$ & $163$\\
$(1,3,0)$ & $(0,3,0)$ & $34$\\
$(1,3,3)$ & $(0,3,3)$ & $39$\\
$(2,3,2)$ & $(1,3,2)$ & $78$\\
$(3,2,2)$ & $(2,2,2)$ & $110$\\
$(3,3,1)$ & $(2,3,1)$ & $116$\\
$(4,2,1)$ & $(3,2,1)$ & $147$\\
$(0,1,2)$ & $(0,0,2)$ & $7$\\
$(4,1,0)$ & $(4,1,1)$ & $167$\\
$(2,3,0)$ & $(1,3,0)$ & $74$\\
$(2,3,3)$ & $(1,3,3)$ & $79$\\
$(3,3,2)$ & $(2,3,2)$ & $118$\\
$(4,3,1)$ & $(4,2,1)$ & $177$\\
$(1,1,2)$ & $(0,1,2)$ & $19$\\
$(4,2,0)$ & $(4,1,0)$ & $168$\\
$(3,3,0)$ & $(2,3,0)$ & $114$\\
$(3,3,3)$ & $(2,3,3)$ & $119$\\
$(4,3,2)$ & $(3,3,2)$ & $158$\\
$(1,1,3)$ & $(1,1,2)$ & $57$\\
$(2,1,2)$ & $(1,1,2)$ & $59$\\
$(4,3,0)$ & $(3,3,0)$ & $154$\\
$(4,3,3)$ & $(3,3,3)$ & $159$\\
$(0,1,3)$ & $(1,1,3)$ & $21$\\
$(2,1,3)$ & $(1,1,3)$ & $61$\\
$(3,1,2)$ & $(2,1,2)$ & $99$\\
$(0,0,3)$ & $(0,1,3)$ & $9$\\
$(2,2,3)$ & $(2,1,3)$ & $100$\\
$(3,1,3)$ & $(2,1,3)$ & $101$\\
$(4,1,2)$ & $(3,1,2)$ & $139$\\
$(1,0,3)$ & $(0,0,3)$ & $10$\\
$(1,2,3)$ & $(2,2,3)$ & $72$\\
$(3,2,3)$ & $(2,2,3)$ & $112$\\
$(4,1,3)$ & $(4,1,2)$ & $171$\\
$(4,2,2)$ & $(4,1,2)$ & $172$\\
$(2,0,3)$ & $(1,0,3)$ & $50$\\
$(0,2,3)$ & $(1,2,3)$ & $32$\\
$(4,2,3)$ & $(3,2,3)$ & $152$\\
$(3,0,3)$ & $(2,0,3)$ & $90$\\
$(4,0,3)$ & $(3,0,3)$ & $130$\\
\bottomrule
\end{longtable}
Thus all 80 vertices are reached by 79 selected edges, proving that the selected graph is a tree.  Removing the 15 retained chords partitions it into the following 16 components; the left entry is the unique terminal in the component.
\begin{longtable}{@{}p{.13\textwidth}p{.79\textwidth}@{}}
\toprule terminal & complete vertex set\\ \midrule
\endfirsthead \toprule terminal & complete vertex set\\ \midrule \endhead
$(4,0,0)$ & $(0,0,0)$, $(0,1,0)$, $(1,0,0)$, $(1,1,0)$, $(2,0,0)$, $(2,1,0)$, $(3,0,0)$, $(3,1,0)$, $(4,0,0)$\\
$(4,0,1)$ & $(0,0,1)$, $(1,0,1)$, $(2,0,1)$, $(3,0,1)$, $(4,0,1)$\\
$(4,0,2)$ & $(4,0,2)$\\
$(4,0,3)$ & $(0,0,3)$, $(0,1,3)$, $(1,0,3)$, $(1,1,3)$, $(2,0,3)$, $(2,1,3)$, $(3,0,3)$, $(3,1,3)$, $(4,0,3)$\\
$(4,1,0)$ & $(4,1,0)$\\
$(4,1,1)$ & $(4,1,1)$\\
$(4,1,2)$ & $(0,0,2)$, $(0,1,2)$, $(1,0,2)$, $(1,1,2)$, $(2,0,2)$, $(2,1,2)$, $(3,0,2)$, $(3,1,2)$, $(4,1,2)$\\
$(4,1,3)$ & $(4,1,3)$\\
$(4,2,0)$ & $(4,2,0)$\\
$(4,2,1)$ & $(0,1,1)$, $(0,2,1)$, $(0,3,1)$, $(1,1,1)$, $(1,2,1)$, $(1,3,1)$, $(2,1,1)$, $(2,2,1)$, $(2,3,1)$, $(3,1,1)$, $(3,2,1)$, $(3,3,1)$, $(4,2,1)$\\
$(4,2,2)$ & $(4,2,2)$\\
$(4,2,3)$ & $(0,2,3)$, $(1,2,3)$, $(2,2,3)$, $(3,2,3)$, $(4,2,3)$\\
$(4,3,0)$ & $(0,2,0)$, $(0,3,0)$, $(1,2,0)$, $(1,3,0)$, $(2,2,0)$, $(2,3,0)$, $(3,2,0)$, $(3,3,0)$, $(4,3,0)$\\
$(4,3,1)$ & $(4,3,1)$\\
$(4,3,2)$ & $(1,3,2)$, $(2,3,2)$, $(3,3,2)$, $(4,3,2)$\\
$(4,3,3)$ & $(0,2,2)$, $(0,3,2)$, $(0,3,3)$, $(1,2,2)$, $(1,3,3)$, $(2,2,2)$, $(2,3,3)$, $(3,2,2)$, $(3,3,3)$, $(4,3,3)$\\
\bottomrule
\end{longtable}
Finally delete from the face dual the duals of the 79 gauge-tree edges and the 15 retained chords.  The remainder has 90 face vertices.  The following 89 parent incidences reach every face from $O$; the last column is the primal grid edge crossed by the dual parent edge.
\begin{longtable}{@{}rrr@{}}
\toprule child face & parent face & primal edge index\\ \midrule
\endfirsthead \toprule child face & parent face & primal edge index\\ \midrule \endhead
$(0;4,1,2)$ & $O$ & $173$\\
$(0;4,2,1)$ & $O$ & $182$\\
$(1;3,0,0)$ & $O$ & $160$\\
$(1;3,0,2)$ & $O$ & $164$\\
$(1;3,3,0)$ & $O$ & $181$\\
$(1;3,3,2)$ & $O$ & $183$\\
$(2;3,0,0)$ & $O$ & $161$\\
$(2;3,0,3)$ & $O$ & $166$\\
$(2;3,2,0)$ & $O$ & $175$\\
$(2;3,2,3)$ & $O$ & $180$\\
$(1;3,2,2)$ & $(0;4,1,2)$ & $178$\\
$(0;4,1,1)$ & $(0;4,2,1)$ & $176$\\
$(2;3,2,2)$ & $(0;4,2,1)$ & $179$\\
$(1;2,0,0)$ & $(1;3,0,0)$ & $120$\\
$(1;2,0,2)$ & $(1;3,0,2)$ & $126$\\
$(1;2,3,0)$ & $(1;3,3,0)$ & $153$\\
$(1;2,3,2)$ & $(1;3,3,2)$ & $157$\\
$(2;2,0,0)$ & $(2;3,0,0)$ & $121$\\
$(2;2,0,3)$ & $(2;3,0,3)$ & $129$\\
$(2;2,2,0)$ & $(2;3,2,0)$ & $143$\\
$(2;2,2,3)$ & $(2;3,2,3)$ & $151$\\
$(0;3,1,2)$ & $(1;3,2,2)$ & $148$\\
$(0;4,0,1)$ & $(0;4,1,1)$ & $169$\\
$(0;4,1,0)$ & $(0;4,1,1)$ & $170$\\
$(0;3,2,1)$ & $(2;3,2,2)$ & $149$\\
$(1;1,0,0)$ & $(1;2,0,0)$ & $80$\\
$(1;1,0,2)$ & $(1;2,0,2)$ & $86$\\
$(1;1,3,0)$ & $(1;2,3,0)$ & $113$\\
$(1;1,3,2)$ & $(1;2,3,2)$ & $117$\\
$(2;1,0,0)$ & $(2;2,0,0)$ & $81$\\
$(2;1,0,3)$ & $(2;2,0,3)$ & $89$\\
$(2;1,2,0)$ & $(2;2,2,0)$ & $103$\\
$(2;1,2,3)$ & $(2;2,2,3)$ & $111$\\
$(1;3,1,2)$ & $(0;3,1,2)$ & $137$\\
$(2;2,1,2)$ & $(0;3,1,2)$ & $138$\\
$(2;2,1,3)$ & $(0;3,1,2)$ & $140$\\
$(2;3,0,2)$ & $(0;4,0,1)$ & $165$\\
$(1;3,2,0)$ & $(0;4,1,0)$ & $174$\\
$(1;2,2,1)$ & $(0;3,2,1)$ & $145$\\
$(1;2,3,1)$ & $(0;3,2,1)$ & $155$\\
$(2;3,2,1)$ & $(0;3,2,1)$ & $146$\\
$(1;0,0,0)$ & $(1;1,0,0)$ & $40$\\
$(1;0,0,2)$ & $(1;1,0,2)$ & $46$\\
$(1;0,3,0)$ & $(1;1,3,0)$ & $73$\\
$(1;0,3,2)$ & $(1;1,3,2)$ & $77$\\
$(2;0,0,0)$ & $(2;1,0,0)$ & $41$\\
$(2;0,0,3)$ & $(2;1,0,3)$ & $49$\\
$(2;0,2,0)$ & $(2;1,2,0)$ & $63$\\
$(2;0,2,3)$ & $(2;1,2,3)$ & $71$\\
$(0;2,1,2)$ & $(2;2,1,2)$ & $98$\\
$(0;3,0,1)$ & $(2;3,0,2)$ & $127$\\
$(0;3,1,0)$ & $(1;3,2,0)$ & $142$\\
$(0;2,2,1)$ & $(1;2,2,1)$ & $105$\\
$(0;0,2,0)$ & $(1;0,3,0)$ & $33$\\
$(0;0,2,2)$ & $(1;0,3,2)$ & $37$\\
$(1;1,1,2)$ & $(0;2,1,2)$ & $97$\\
$(1;1,2,2)$ & $(0;2,1,2)$ & $108$\\
$(1;2,0,1)$ & $(0;3,0,1)$ & $123$\\
$(1;2,1,1)$ & $(0;3,0,1)$ & $134$\\
$(2;3,0,1)$ & $(0;3,0,1)$ & $124$\\
$(1;3,1,0)$ & $(0;3,1,0)$ & $131$\\
$(2;2,1,0)$ & $(0;3,1,0)$ & $132$\\
$(2;2,1,1)$ & $(0;3,1,0)$ & $135$\\
$(2;1,2,1)$ & $(0;2,2,1)$ & $106$\\
$(2;1,2,2)$ & $(0;2,2,1)$ & $109$\\
$(0;0,1,0)$ & $(0;0,2,0)$ & $22$\\
$(0;0,1,2)$ & $(0;0,2,2)$ & $28$\\
$(0;1,1,2)$ & $(1;1,2,2)$ & $68$\\
$(0;2,0,1)$ & $(1;2,1,1)$ & $94$\\
$(0;2,1,0)$ & $(2;2,1,1)$ & $95$\\
$(0;1,2,1)$ & $(2;1,2,1)$ & $66$\\
$(0;0,0,0)$ & $(0;0,1,0)$ & $11$\\
$(0;0,0,2)$ & $(0;0,1,2)$ & $17$\\
$(2;0,1,2)$ & $(0;1,1,2)$ & $58$\\
$(2;0,1,3)$ & $(0;1,1,2)$ & $60$\\
$(2;1,0,1)$ & $(0;2,0,1)$ & $84$\\
$(2;1,0,2)$ & $(0;2,0,1)$ & $87$\\
$(1;1,1,0)$ & $(0;2,1,0)$ & $91$\\
$(1;1,2,0)$ & $(0;2,1,0)$ & $102$\\
$(1;0,2,1)$ & $(0;1,2,1)$ & $65$\\
$(1;0,3,1)$ & $(0;1,2,1)$ & $75$\\
$(0;1,0,1)$ & $(2;1,0,2)$ & $47$\\
$(0;1,1,0)$ & $(1;1,2,0)$ & $62$\\
$(0;0,2,1)$ & $(1;0,3,1)$ & $35$\\
$(1;0,0,1)$ & $(0;1,0,1)$ & $43$\\
$(1;0,1,1)$ & $(0;1,0,1)$ & $54$\\
$(2;0,1,0)$ & $(0;1,1,0)$ & $52$\\
$(2;0,1,1)$ & $(0;1,1,0)$ & $55$\\
$(0;0,0,1)$ & $(1;0,1,1)$ & $14$\\
\bottomrule
\end{longtable}
The table is a spanning tree of the retained dual remainder and therefore proves its connectivity.  Together, these three printed witnesses establish the tree, one-terminal-per-component, and dual connectivity assertions for the normal width-four base.

\section{Complete opposite width-four encoder base trace}
\label{app:width4-opposite-trace}
This appendix discharges the finite base case used by the opposite shell induction.  Vertices and edge indices use the convention of \cref{app:encoders}.  A face $(a;x,y,z)$ is the unit square with fixed coordinate axis $a$ and lower corner $(x,y,z)$; $O$ denotes the nonsquare outer face.  Each row below is an explicit incidence check, so no program output is needed to infer connectivity.
The opposite base contains 80 vertices and the 79 selected tree edges printed in \cref{app:encoders}.  Starting at $(0,0,0)$, the following table gives one parent edge for every other vertex.
\begin{longtable}{@{}rrr@{}}
\toprule child & parent & edge index\\ \midrule
\endfirsthead \toprule child & parent & edge index\\ \midrule \endhead
$(0,1,0)$ & $(0,0,0)$ & $1$\\
$(1,0,0)$ & $(0,0,0)$ & $2$\\
$(0,1,1)$ & $(0,1,0)$ & $11$\\
$(2,0,0)$ & $(1,0,0)$ & $42$\\
$(0,0,1)$ & $(0,1,1)$ & $4$\\
$(0,2,1)$ & $(0,1,1)$ & $15$\\
$(1,1,1)$ & $(0,1,1)$ & $16$\\
$(3,0,0)$ & $(2,0,0)$ & $82$\\
$(1,0,1)$ & $(0,0,1)$ & $5$\\
$(0,3,1)$ & $(0,2,1)$ & $26$\\
$(1,2,1)$ & $(0,2,1)$ & $27$\\
$(2,1,1)$ & $(1,1,1)$ & $56$\\
$(4,0,0)$ & $(3,0,0)$ & $122$\\
$(1,0,2)$ & $(1,0,1)$ & $43$\\
$(2,0,1)$ & $(1,0,1)$ & $45$\\
$(1,3,1)$ & $(0,3,1)$ & $36$\\
$(1,2,2)$ & $(1,2,1)$ & $65$\\
$(2,2,1)$ & $(1,2,1)$ & $67$\\
$(2,1,0)$ & $(2,1,1)$ & $91$\\
$(3,1,1)$ & $(2,1,1)$ & $96$\\
$(0,0,2)$ & $(1,0,2)$ & $8$\\
$(2,0,2)$ & $(1,0,2)$ & $48$\\
$(3,0,1)$ & $(2,0,1)$ & $85$\\
$(2,3,1)$ & $(1,3,1)$ & $76$\\
$(0,2,2)$ & $(1,2,2)$ & $30$\\
$(2,2,2)$ & $(1,2,2)$ & $70$\\
$(3,2,1)$ & $(2,2,1)$ & $107$\\
$(1,1,0)$ & $(2,1,0)$ & $53$\\
$(3,1,0)$ & $(2,1,0)$ & $93$\\
$(0,1,2)$ & $(0,0,2)$ & $7$\\
$(3,0,2)$ & $(2,0,2)$ & $88$\\
$(4,0,1)$ & $(3,0,1)$ & $125$\\
$(3,3,1)$ & $(2,3,1)$ & $116$\\
$(0,3,2)$ & $(0,2,2)$ & $29$\\
$(2,3,2)$ & $(2,2,2)$ & $109$\\
$(3,2,2)$ & $(2,2,2)$ & $110$\\
$(1,2,0)$ & $(1,1,0)$ & $52$\\
$(4,1,0)$ & $(3,1,0)$ & $133$\\
$(0,1,3)$ & $(0,1,2)$ & $17$\\
$(1,1,2)$ & $(0,1,2)$ & $19$\\
$(4,0,2)$ & $(4,0,1)$ & $162$\\
$(4,1,1)$ & $(4,0,1)$ & $163$\\
$(0,3,3)$ & $(0,3,2)$ & $37$\\
$(1,3,2)$ & $(2,3,2)$ & $78$\\
$(3,3,2)$ & $(2,3,2)$ & $118$\\
$(0,2,0)$ & $(1,2,0)$ & $24$\\
$(2,2,0)$ & $(1,2,0)$ & $64$\\
$(4,2,0)$ & $(4,1,0)$ & $168$\\
$(0,0,3)$ & $(0,1,3)$ & $9$\\
$(2,1,2)$ & $(1,1,2)$ & $59$\\
$(1,3,3)$ & $(0,3,3)$ & $39$\\
$(4,3,2)$ & $(3,3,2)$ & $158$\\
$(0,3,0)$ & $(0,2,0)$ & $23$\\
$(3,2,0)$ & $(2,2,0)$ & $104$\\
$(1,0,3)$ & $(0,0,3)$ & $10$\\
$(2,1,3)$ & $(2,1,2)$ & $97$\\
$(3,1,2)$ & $(2,1,2)$ & $99$\\
$(2,3,3)$ & $(1,3,3)$ & $79$\\
$(1,3,0)$ & $(0,3,0)$ & $34$\\
$(2,0,3)$ & $(1,0,3)$ & $50$\\
$(1,1,3)$ & $(2,1,3)$ & $61$\\
$(3,1,3)$ & $(2,1,3)$ & $101$\\
$(4,1,2)$ & $(3,1,2)$ & $139$\\
$(3,3,3)$ & $(2,3,3)$ & $119$\\
$(2,3,0)$ & $(1,3,0)$ & $74$\\
$(3,0,3)$ & $(2,0,3)$ & $90$\\
$(1,2,3)$ & $(1,1,3)$ & $60$\\
$(4,1,3)$ & $(3,1,3)$ & $141$\\
$(4,2,2)$ & $(4,1,2)$ & $172$\\
$(4,3,3)$ & $(3,3,3)$ & $159$\\
$(3,3,0)$ & $(2,3,0)$ & $114$\\
$(4,0,3)$ & $(3,0,3)$ & $130$\\
$(0,2,3)$ & $(1,2,3)$ & $32$\\
$(2,2,3)$ & $(1,2,3)$ & $72$\\
$(4,2,1)$ & $(4,2,2)$ & $176$\\
$(4,3,0)$ & $(3,3,0)$ & $154$\\
$(3,2,3)$ & $(2,2,3)$ & $112$\\
$(4,3,1)$ & $(4,2,1)$ & $177$\\
$(4,2,3)$ & $(3,2,3)$ & $152$\\
\bottomrule
\end{longtable}
Thus all 80 vertices are reached by 79 selected edges, proving that the selected graph is a tree.  Removing the 15 retained chords partitions it into the following 16 components; the left entry is the unique terminal in the component.
\begin{longtable}{@{}p{.13\textwidth}p{.79\textwidth}@{}}
\toprule terminal & complete vertex set\\ \midrule
\endfirsthead \toprule terminal & complete vertex set\\ \midrule \endhead
$(4,0,0)$ & $(0,0,0)$, $(0,1,0)$, $(1,0,0)$, $(2,0,0)$, $(3,0,0)$, $(4,0,0)$\\
$(4,0,1)$ & $(0,0,1)$, $(0,1,1)$, $(0,2,1)$, $(0,3,1)$, $(1,0,1)$, $(1,1,1)$, $(1,2,1)$, $(1,3,1)$, $(2,0,1)$, $(2,1,1)$, $(2,2,1)$, $(2,3,1)$, $(3,0,1)$, $(3,1,1)$, $(3,2,1)$, $(3,3,1)$, $(4,0,1)$\\
$(4,0,2)$ & $(4,0,2)$\\
$(4,0,3)$ & $(0,0,3)$, $(0,1,3)$, $(1,0,3)$, $(2,0,3)$, $(3,0,3)$, $(4,0,3)$\\
$(4,1,0)$ & $(1,1,0)$, $(2,1,0)$, $(3,1,0)$, $(4,1,0)$\\
$(4,1,1)$ & $(4,1,1)$\\
$(4,1,2)$ & $(0,0,2)$, $(0,1,2)$, $(1,0,2)$, $(1,1,2)$, $(2,0,2)$, $(2,1,2)$, $(3,0,2)$, $(3,1,2)$, $(4,1,2)$\\
$(4,1,3)$ & $(1,1,3)$, $(2,1,3)$, $(3,1,3)$, $(4,1,3)$\\
$(4,2,0)$ & $(4,2,0)$\\
$(4,2,1)$ & $(4,2,1)$\\
$(4,2,2)$ & $(4,2,2)$\\
$(4,2,3)$ & $(0,2,3)$, $(1,2,3)$, $(2,2,3)$, $(3,2,3)$, $(4,2,3)$\\
$(4,3,0)$ & $(0,2,0)$, $(0,3,0)$, $(1,2,0)$, $(1,3,0)$, $(2,2,0)$, $(2,3,0)$, $(3,2,0)$, $(3,3,0)$, $(4,3,0)$\\
$(4,3,1)$ & $(4,3,1)$\\
$(4,3,2)$ & $(1,3,2)$, $(2,3,2)$, $(3,3,2)$, $(4,3,2)$\\
$(4,3,3)$ & $(0,2,2)$, $(0,3,2)$, $(0,3,3)$, $(1,2,2)$, $(1,3,3)$, $(2,2,2)$, $(2,3,3)$, $(3,2,2)$, $(3,3,3)$, $(4,3,3)$\\
\bottomrule
\end{longtable}
Finally delete from the face dual the duals of the 79 gauge-tree edges and the 15 retained chords.  The remainder has 90 face vertices.  The following 89 parent incidences reach every face from $O$; the last column is the primal grid edge crossed by the dual parent edge.
\begin{longtable}{@{}rrr@{}}
\toprule child face & parent face & primal edge index\\ \midrule
\endfirsthead \toprule child face & parent face & primal edge index\\ \midrule \endhead
$(0;0,2,1)$ & $O$ & $35$\\
$(1;0,3,0)$ & $O$ & $33$\\
$(1;0,3,2)$ & $O$ & $37$\\
$(0;0,1,1)$ & $(0;0,2,1)$ & $25$\\
$(1;1,3,0)$ & $(1;0,3,0)$ & $73$\\
$(1;1,3,2)$ & $(1;0,3,2)$ & $77$\\
$(0;0,0,1)$ & $(0;0,1,1)$ & $14$\\
$(1;2,3,0)$ & $(1;1,3,0)$ & $113$\\
$(1;2,3,2)$ & $(1;1,3,2)$ & $117$\\
$(1;3,3,0)$ & $(1;2,3,0)$ & $153$\\
$(1;3,3,2)$ & $(1;2,3,2)$ & $157$\\
$(0;4,2,0)$ & $(1;3,3,0)$ & $181$\\
$(0;4,2,2)$ & $(1;3,3,2)$ & $183$\\
$(0;4,1,0)$ & $(0;4,2,0)$ & $174$\\
$(2;3,2,0)$ & $(0;4,2,0)$ & $175$\\
$(0;4,1,2)$ & $(0;4,2,2)$ & $178$\\
$(0;4,2,1)$ & $(0;4,2,2)$ & $179$\\
$(2;3,2,3)$ & $(0;4,2,2)$ & $180$\\
$(0;4,0,0)$ & $(0;4,1,0)$ & $167$\\
$(2;3,1,1)$ & $(0;4,1,0)$ & $170$\\
$(2;2,2,0)$ & $(2;3,2,0)$ & $143$\\
$(0;4,0,2)$ & $(0;4,1,2)$ & $171$\\
$(2;3,1,3)$ & $(0;4,1,2)$ & $173$\\
$(1;3,3,1)$ & $(0;4,2,1)$ & $182$\\
$(2;2,2,3)$ & $(2;3,2,3)$ & $151$\\
$(1;3,0,0)$ & $(0;4,0,0)$ & $160$\\
$(2;3,0,0)$ & $(0;4,0,0)$ & $161$\\
$(0;3,1,0)$ & $(2;3,1,1)$ & $135$\\
$(2;1,2,0)$ & $(2;2,2,0)$ & $103$\\
$(0;4,0,1)$ & $(0;4,0,2)$ & $165$\\
$(1;3,0,2)$ & $(0;4,0,2)$ & $164$\\
$(2;3,0,3)$ & $(0;4,0,2)$ & $166$\\
$(0;3,1,2)$ & $(2;3,1,3)$ & $140$\\
$(0;3,2,1)$ & $(1;3,3,1)$ & $155$\\
$(2;1,2,3)$ & $(2;2,2,3)$ & $111$\\
$(1;2,0,0)$ & $(1;3,0,0)$ & $120$\\
$(2;2,0,0)$ & $(2;3,0,0)$ & $121$\\
$(1;2,1,0)$ & $(0;3,1,0)$ & $131$\\
$(1;2,2,0)$ & $(0;3,1,0)$ & $142$\\
$(2;3,1,0)$ & $(0;3,1,0)$ & $132$\\
$(2;0,2,0)$ & $(2;1,2,0)$ & $63$\\
$(1;3,1,1)$ & $(0;4,0,1)$ & $169$\\
$(1;2,0,2)$ & $(1;3,0,2)$ & $126$\\
$(2;2,0,3)$ & $(2;3,0,3)$ & $129$\\
$(1;2,1,2)$ & $(0;3,1,2)$ & $137$\\
$(1;2,2,2)$ & $(0;3,1,2)$ & $148$\\
$(2;3,1,2)$ & $(0;3,1,2)$ & $138$\\
$(1;3,2,1)$ & $(0;3,2,1)$ & $145$\\
$(2;2,2,1)$ & $(0;3,2,1)$ & $146$\\
$(2;2,2,2)$ & $(0;3,2,1)$ & $149$\\
$(2;0,2,3)$ & $(2;1,2,3)$ & $71$\\
$(1;1,0,0)$ & $(1;2,0,0)$ & $80$\\
$(2;1,0,0)$ & $(2;2,0,0)$ & $81$\\
$(0;2,1,0)$ & $(1;2,2,0)$ & $102$\\
$(0;3,0,1)$ & $(1;3,1,1)$ & $134$\\
$(1;1,0,2)$ & $(1;2,0,2)$ & $86$\\
$(2;1,0,3)$ & $(2;2,0,3)$ & $89$\\
$(0;2,1,2)$ & $(1;2,2,2)$ & $108$\\
$(0;2,2,1)$ & $(2;2,2,1)$ & $106$\\
$(1;0,0,0)$ & $(1;1,0,0)$ & $40$\\
$(2;0,0,0)$ & $(2;1,0,0)$ & $41$\\
$(2;1,1,0)$ & $(0;2,1,0)$ & $92$\\
$(2;1,1,1)$ & $(0;2,1,0)$ & $95$\\
$(1;3,0,1)$ & $(0;3,0,1)$ & $123$\\
$(2;2,0,1)$ & $(0;3,0,1)$ & $124$\\
$(2;2,0,2)$ & $(0;3,0,1)$ & $127$\\
$(1;0,0,2)$ & $(1;1,0,2)$ & $46$\\
$(2;0,0,3)$ & $(2;1,0,3)$ & $49$\\
$(2;1,1,2)$ & $(0;2,1,2)$ & $98$\\
$(2;1,1,3)$ & $(0;2,1,2)$ & $100$\\
$(1;1,2,1)$ & $(0;2,2,1)$ & $105$\\
$(1;1,3,1)$ & $(0;2,2,1)$ & $115$\\
$(0;1,1,0)$ & $(2;1,1,1)$ & $55$\\
$(0;2,0,1)$ & $(2;2,0,1)$ & $84$\\
$(0;1,1,2)$ & $(2;1,1,2)$ & $58$\\
$(0;1,2,1)$ & $(1;1,3,1)$ & $75$\\
$(1;0,1,0)$ & $(0;1,1,0)$ & $51$\\
$(1;0,2,0)$ & $(0;1,1,0)$ & $62$\\
$(1;1,0,1)$ & $(0;2,0,1)$ & $83$\\
$(1;1,1,1)$ & $(0;2,0,1)$ & $94$\\
$(1;0,1,2)$ & $(0;1,1,2)$ & $57$\\
$(1;0,2,2)$ & $(0;1,1,2)$ & $68$\\
$(2;0,2,1)$ & $(0;1,2,1)$ & $66$\\
$(2;0,2,2)$ & $(0;1,2,1)$ & $69$\\
$(0;0,1,0)$ & $(1;0,2,0)$ & $22$\\
$(0;1,0,1)$ & $(1;1,1,1)$ & $54$\\
$(0;0,1,2)$ & $(1;0,2,2)$ & $28$\\
$(2;0,0,1)$ & $(0;1,0,1)$ & $44$\\
$(2;0,0,2)$ & $(0;1,0,1)$ & $47$\\
\bottomrule
\end{longtable}
The table is a spanning tree of the retained dual remainder and therefore proves its connectivity.  Together, these three printed witnesses establish the tree, one-terminal-per-component, and dual connectivity assertions for the opposite width-four base.
For the exceptional width-four relation, remove also edge 37 from the retained dual remainder.  The component containing $O$ is the following complete face set:
\begin{quote}\small $O$, $(1;0,0,0)$, $(2;0,0,0)$, $(0;0,0,1)$, $(1;0,1,0)$, $(0;0,1,0)$, $(0;0,1,1)$, $(1;0,2,0)$, $(2;0,2,0)$, $(0;0,2,1)$, $(1;0,3,0)$, $(1;1,0,0)$, $(2;1,0,0)$, $(2;1,1,0)$, $(0;1,1,0)$, $(2;1,1,1)$, $(2;1,2,0)$, $(1;1,3,0)$, $(1;2,0,0)$, $(2;2,0,0)$, $(1;2,1,0)$, $(0;2,1,0)$, $(1;2,2,0)$, $(2;2,2,0)$, $(1;2,3,0)$, $(1;3,0,0)$, $(2;3,0,0)$, $(2;3,1,0)$, $(0;3,1,0)$, $(2;3,1,1)$, $(2;3,2,0)$, $(1;3,3,0)$, $(0;4,0,0)$, $(0;4,1,0)$, $(0;4,2,0)$.\end{quote}
Its non-gauge chord boundary consists exactly of primal edge indices $37,163,177$, namely $X=e_z(0,3,2)$, $e_y(4,0,1)$, and $e_y(4,2,1)$.  Its face-boundary relation is therefore $h_X=h_{e_y(4,0,1)}+h_{e_y(4,2,1)}$.
For the terminal relation, deleting edge 37 leaves terminal set $\{(4,3,3)\}$ away from the root; deleting edge 65 ($e_z(1,2,1)$) leaves $\{(4,3,2),(4,3,3)\}$; and deleting edge 109 ($e_y(2,2,2)$) leaves $\{(4,3,2)\}$.  Symmetric difference therefore gives $u_X=u_{e_z(1,2,1)}+u_{e_y(2,2,2)}$.


\clearpage
\section{Reproducibility and claim inventory}

Zenodo version DOI \url{https://doi.org/10.5281/zenodo.21845273} is reserved
in an unsubmitted draft record.  The DOI is not yet registered, no files have
been uploaded, and there is no public immutable proof archive; the manuscript
is therefore not submission-ready.  Replay links will be inserted only when
they identify the exact archived revision.

The proof archive separates discovery scripts from exact proof scripts.  The
authoritative immutable records are:
\begin{center}
\small
\begin{tabular}{@{}ll@{}}
\toprule
claim & artifact\\
\midrule
grid upper bound & \texttt{cycle-7-b7-lane-b-arbitrary-width-closure-v1.json}\\
generic tightness & \texttt{cycle-8-b11-g1-generic-tightness-v2.json}\\
Walsh marginals & \texttt{cycle-9-b8-lane-b-all-q-marginals-v1.json}\\
abstract theorem and \(K_{3,3}\) chain &
\texttt{cycle-10-b9-abstract-separator-k33-sharpness-v1.json}\\
embedding robustness &
\texttt{cycle-11-b10-lane-b-embedding-robustness-v2.json}\\
\bottomrule
\end{tabular}
\end{center}
Each record freezes source hashes, conventions, prime specializations, and a
one-command immutable check.  The corrected generic-tightness record removes
volatile timing fields from its predecessor; no mathematical field changed.
From the project directory the principal checks are
\begin{verbatim}
python3 proof/build_cycle7_lane_b_arbitrary_width_closure.py --check
python3 proof/build_cycle8_g1_generic_tightness_correction_v2.py --check
python3 proof/build_cycle10_abstract_separator_theorem.py --check
python3 discovery/diagnose_g1_opposite_quotient.py --maximum-width 20
python3 discovery/render_width4_encoder_base_trace.py --phase normal
python3 discovery/render_width4_encoder_base_trace.py --phase opposite
\end{verbatim}
The first three commands audit frozen claims.  The last three regenerate the
working quotient and printed base tables.  None replaces the \(\forall w\)
arguments in the manuscript.

The principal proof documents are:
\begin{itemize}
\item \path{proof/lane_b_arbitrary_width_closure_proof.md};
\item \path{proof/g1_arbitrary_width_generic_tightness.md};
\item \path{proof/abstract_spin_structure_separator_theorem.md};
\item \path{proof/lane_b_all_q_marginal_algorithm.md}.
\end{itemize}
The replay suite
pins Python \(3.11\) and uses exact integer, rational, and finite-field
arithmetic.  Important modular certificates are replayed over the two primes
listed in \cref{tab:ranks}; denominators and normalization factors are checked
invertible before reduction.

\paragraph{AI assistance.}
The author used AI assistance in developing candidate formulations, proof
presentation, and verification code.  Claims marked \proofpending are not
promoted mathematical results; exact replays are identified separately as
audits.

\bibliographystyle{alpha}
\bibliography{references}

\end{document}
